\documentclass{otscience}
\def\OTCOMBINED{1}
% ============================================================
% Shared setup for OT Science modular workbook parts
% Place these files in the same folder as your fixed otscience.cls
% ============================================================

\makeatletter
\makeatother
\usepackage{multicol}
\usepackage{longtable}
\usepackage{makecell}
\usepackage{pdflscape}
\usepackage{bm}
\providecommand{\blank}[1][3cm]{\rule{#1}{0.4pt}}
\providecommand{\bigblank}{\rule{7cm}{0.4pt}}
\providecommand{\componentblank}{\blank[2.5cm]}
\providecommand{\R}{\mathbb{R}}
\providecommand{\C}{\mathbb{C}}
\providecommand{\ii}{\mathrm{i}}
\providecommand{\ee}{\mathrm{e}}
\providecommand{\dd}{\mathrm{d}}
\providecommand{\grad}{\nabla}
\providecommand{\divergence}{\nabla\!\cdot}
\providecommand{\curl}{\nabla\!\times}
\providecommand{\lap}{\nabla^2}
\providecommand{\ihat}{\hat{\mathbf{i}}}
\providecommand{\jhat}{\hat{\mathbf{j}}}
\providecommand{\khat}{\hat{\mathbf{k}}}
\providecommand{\er}{\hat{\mathbf{e}}_r}
\providecommand{\etheta}{\hat{\mathbf{e}}_\theta}
\providecommand{\ephi}{\hat{\mathbf{e}}_\phi}
\providecommand{\erho}{\hat{\boldsymbol{\rho}}}
\providecommand{\vA}{\mathbf{A}}
\providecommand{\vB}{\mathbf{B}}
\providecommand{\va}{\mathbf{a}}
\providecommand{\vb}{\mathbf{b}}
\providecommand{\vc}{\mathbf{c}}
\providecommand{\vx}{\mathbf{x}}
\providecommand{\vr}{\mathbf{r}}
\providecommand{\matA}{\mathbf{A}}
\providecommand{\matB}{\mathbf{B}}
\providecommand{\tr}{\operatorname{tr}}
\providecommand{\cof}{\operatorname{cof}}
\providecommand{\dev}{\operatorname{dev}}
\providecommand{\diag}{\operatorname{diag}}
\providecommand{\questionline}[1][5cm]{\rule{#1}{0.4pt}}
\setlength{\parindent}{0pt}
\setlength{\parskip}{4pt}
\renewcommand{\arraystretch}{1.25}

% Fallbacks in case the current otscience.cls has not defined nosplit boxes.
\makeatletter
\@ifundefined{formulaboxnosplit}{\newenvironment{formulaboxnosplit}[1][Formula]{\begin{formulabox}[#1]}{\end{formulabox}}}{}
\@ifundefined{definitionboxnosplit}{\newenvironment{definitionboxnosplit}[1][Definition]{\begin{definitionbox}[#1]}{\end{definitionbox}}}{}
\@ifundefined{summaryboxnosplit}{\newenvironment{summaryboxnosplit}[1][Summary]{\begin{summarybox}[#1]}{\end{summarybox}}}{}
\@ifundefined{warningboxnosplit}{\newenvironment{warningboxnosplit}[1][Warning]{\begin{warningbox}[#1]}{\end{warningbox}}}{}
\@ifundefined{exampleboxnosplit}{\newenvironment{exampleboxnosplit}[1][Example]{\begin{examplebox}[#1]}{\end{examplebox}}}{}
\makeatother

\newenvironment{practicebox}[1][Quick Drills and Practice]{\begin{examplebox}[#1]}{\end{examplebox}}
\newenvironment{practiceboxnosplit}[1][Quick Drills and Practice]{\begin{exampleboxnosplit}[#1]}{\end{exampleboxnosplit}}

\newcommand{\standalonetitle}[3]{%
    \title{\textbf{#1}\\\large #2}
    \author{OT Science Notes}
    \date{#3}
}

\usepackage{silence}

\standalonetitle{Vector Calculus, Coordinate Systems, Tensor Operations, and Complex Numbers}{Modular Formula Completion and Practice Workbook}{\DTMtoday}

\begin{document}
\maketitle
\tableofcontents
\newpage

\section*{How to Use the Modular Workbook}
\addcontentsline{toc}{section}{How to Use the Modular Workbook}
This combined file imports each modular part. Each numbered file can also be compiled and printed independently. The design keeps formulas in formula boxes, places practice work in one larger box per major section, and leaves ordinary explanation as normal text so that the page does not become visually noisy.

When a box uses a \texttt{..boxnosplit} environment, it is because the material should remain together on one page where possible. Longer practice boxes are allowed to break across pages.

\ifdefined\OTCOMBINED%
\else
\documentclass{otscience}
% ============================================================
% Shared setup for OT Science modular workbook parts
% Place these files in the same folder as your fixed otscience.cls
% ============================================================

\makeatletter
\makeatother
\usepackage{multicol}
\usepackage{longtable}
\usepackage{makecell}
\usepackage{pdflscape}
\usepackage{bm}
\providecommand{\blank}[1][3cm]{\rule{#1}{0.4pt}}
\providecommand{\bigblank}{\rule{7cm}{0.4pt}}
\providecommand{\componentblank}{\blank[2.5cm]}
\providecommand{\R}{\mathbb{R}}
\providecommand{\C}{\mathbb{C}}
\providecommand{\ii}{\mathrm{i}}
\providecommand{\ee}{\mathrm{e}}
\providecommand{\dd}{\mathrm{d}}
\providecommand{\grad}{\nabla}
\providecommand{\divergence}{\nabla\!\cdot}
\providecommand{\curl}{\nabla\!\times}
\providecommand{\lap}{\nabla^2}
\providecommand{\ihat}{\hat{\mathbf{i}}}
\providecommand{\jhat}{\hat{\mathbf{j}}}
\providecommand{\khat}{\hat{\mathbf{k}}}
\providecommand{\er}{\hat{\mathbf{e}}_r}
\providecommand{\etheta}{\hat{\mathbf{e}}_\theta}
\providecommand{\ephi}{\hat{\mathbf{e}}_\phi}
\providecommand{\erho}{\hat{\boldsymbol{\rho}}}
\providecommand{\vA}{\mathbf{A}}
\providecommand{\vB}{\mathbf{B}}
\providecommand{\va}{\mathbf{a}}
\providecommand{\vb}{\mathbf{b}}
\providecommand{\vc}{\mathbf{c}}
\providecommand{\vx}{\mathbf{x}}
\providecommand{\vr}{\mathbf{r}}
\providecommand{\matA}{\mathbf{A}}
\providecommand{\matB}{\mathbf{B}}
\providecommand{\tr}{\operatorname{tr}}
\providecommand{\cof}{\operatorname{cof}}
\providecommand{\dev}{\operatorname{dev}}
\providecommand{\diag}{\operatorname{diag}}
\providecommand{\questionline}[1][5cm]{\rule{#1}{0.4pt}}
\setlength{\parindent}{0pt}
\setlength{\parskip}{4pt}
\renewcommand{\arraystretch}{1.25}

% Fallbacks in case the current otscience.cls has not defined nosplit boxes.
\makeatletter
\@ifundefined{formulaboxnosplit}{\newenvironment{formulaboxnosplit}[1][Formula]{\begin{formulabox}[#1]}{\end{formulabox}}}{}
\@ifundefined{definitionboxnosplit}{\newenvironment{definitionboxnosplit}[1][Definition]{\begin{definitionbox}[#1]}{\end{definitionbox}}}{}
\@ifundefined{summaryboxnosplit}{\newenvironment{summaryboxnosplit}[1][Summary]{\begin{summarybox}[#1]}{\end{summarybox}}}{}
\@ifundefined{warningboxnosplit}{\newenvironment{warningboxnosplit}[1][Warning]{\begin{warningbox}[#1]}{\end{warningbox}}}{}
\@ifundefined{exampleboxnosplit}{\newenvironment{exampleboxnosplit}[1][Example]{\begin{examplebox}[#1]}{\end{examplebox}}}{}
\makeatother

\newenvironment{practicebox}[1][Quick Drills and Practice]{\begin{examplebox}[#1]}{\end{examplebox}}
\newenvironment{practiceboxnosplit}[1][Quick Drills and Practice]{\begin{exampleboxnosplit}[#1]}{\end{exampleboxnosplit}}

\newcommand{\standalonetitle}[3]{%
    \title{\textbf{#1}\\\large #2}
    \author{OT Science Notes}
    \date{#3}
}

\standalonetitle{Gradient, Divergence and Curl of Vector Fields}{Part 1 of the OT Science Formula Completion Workbook}{\DTMtoday}
\begin{document}
\maketitle
\tableofcontents
\newpage
\fi

\section{Gradient, Divergence and Curl of Vector Fields}

This part fixes the central grammar of vector calculus. The same symbol \(\nabla\) behaves differently depending on what it acts on and how it is combined with dot, cross, or ordinary multiplication. The fastest way to avoid confusion is to always ask what kind of object must come out of the operation.

\begin{center}
    \begin{tabularx}{\textwidth}{>{\bfseries}l X l}
        \toprule
        Operation                  & Input        & Output                               \\
        \midrule
        \(\nabla f\)               & scalar field & vector field                         \\
        \(\nabla\cdot\mathbf{A}\)  & vector field & scalar field                         \\
        \(\nabla\times\mathbf{A}\) & vector field & vector field                         \\
        \(\nabla\mathbf{A}\)       & vector field & second-order tensor/Jacobian         \\
        \(\nabla^2 f\)             & scalar field & scalar field                         \\
        \(\nabla^2\mathbf{A}\)     & vector field & vector field, component by component \\
        \bottomrule
    \end{tabularx}
\end{center}

\subsection{The Cartesian Nabla Operator}

In rectangular Cartesian coordinates, the basis vectors \(\ihat,\jhat,\khat\) are constant. This is why the Cartesian formulas look cleaner than cylindrical or spherical formulas. The derivatives act only on the components of the fields, not on the basis vectors.

\begin{formulaboxnosplit}[Core Cartesian Formulae]
    For a scalar field \(f(x,y,z)\) and a vector field
    \[
        \mathbf{A}=A_x\ihat+A_y\jhat+A_z\khat,
    \]
    \[
        \nabla=\ihat\frac{\partial}{\partial x}+\jhat\frac{\partial}{\partial y}+\khat\frac{\partial}{\partial z}.
    \]

    \[
        \nabla f=\frac{\partial f}{\partial x}\ihat+\frac{\partial f}{\partial y}\jhat+\frac{\partial f}{\partial z}\khat.
    \]

    \[
        \nabla\cdot\mathbf{A}=\frac{\partial A_x}{\partial x}+\frac{\partial A_y}{\partial y}+\frac{\partial A_z}{\partial z}.
    \]

    \[
        \nabla\times\mathbf{A}=\begin{vmatrix}
            \ihat      & \jhat      & \khat      \\
            \partial_x & \partial_y & \partial_z \\
            A_x        & A_y        & A_z
        \end{vmatrix}
    \]

    \[
        =\left(\frac{\partial A_z}{\partial y}-\frac{\partial A_y}{\partial z}\right)\ihat
        +\left(\frac{\partial A_x}{\partial z}-\frac{\partial A_z}{\partial x}\right)\jhat
        +\left(\frac{\partial A_y}{\partial x}-\frac{\partial A_x}{\partial y}\right)\khat.
    \]

    \[
        \nabla^2 f=\frac{\partial^2 f}{\partial x^2}+\frac{\partial^2 f}{\partial y^2}+\frac{\partial^2 f}{\partial z^2}.
    \]
\end{formulaboxnosplit}

\subsection{Gradient of a Scalar Field}

The gradient points in the direction of steepest increase of a scalar field. Its magnitude is the maximum possible directional rate of change. If \(f\) is temperature, then \(\nabla f\) points toward the direction in which temperature increases fastest.

\begin{exampleboxnosplit}[Worked Example: Gradient of a Scalar Field]
    Let
    \[
        f=x^2y+yz^3.
    \]
    Then
    \[
        \frac{\partial f}{\partial x}=2xy,\qquad
        \frac{\partial f}{\partial y}=x^2+z^3,\qquad
        \frac{\partial f}{\partial z}=3yz^2.
    \]
    Therefore
    \[
        \nabla f=2xy\ihat+(x^2+z^3)\jhat+3yz^2\khat.
    \]
\end{exampleboxnosplit}

\subsection{Gradient of a Vector Field}

The expression \(\nabla\mathbf{A}\) is not the same kind of object as \(\nabla f\). The gradient of a scalar is a vector, but the gradient of a vector is a second-order tensor. It tells us how every component of \(\mathbf{A}\) changes in every coordinate direction.

Using the convention
\[
    {(\nabla\mathbf{A})}_{ij}=\frac{\partial A_i}{\partial x_j},
\]
the vector gradient is the matrix
\[
    \nabla\mathbf{A}=
    \begin{pmatrix}
        \partial_x A_x & \partial_y A_x & \partial_z A_x \\
        \partial_x A_y & \partial_y A_y & \partial_z A_y \\
        \partial_x A_z & \partial_y A_z & \partial_z A_z
    \end{pmatrix}.
\]

The divergence is the trace of this tensor:
\[
    \nabla\cdot\mathbf{A}=\tr(\nabla\mathbf{A}).
\]
The curl is related to the skew-symmetric part of \(\nabla\mathbf{A}\). This is why curl is naturally connected to rotation.

\begin{formulabox}[Vector Gradient, Divergence and Curl Relationship]
    \[
        {(\nabla\mathbf{A})}_{ij}=\partial_j A_i.
    \]

    \[
        \nabla\cdot\mathbf{A}=\partial_i A_i=\tr(\nabla\mathbf{A}).
    \]

    \[
        {(\nabla\times\mathbf{A})}_i=\varepsilon_{ijk}\partial_j A_k.
    \]

    If
    \[
        W_{ij}=\frac12(\partial_{j}A_i-\partial_{i}A_j),
    \]
    then \(W_{ij}\) is the skew part of the vector gradient. In three dimensions, the curl is the axial vector associated with this skew part.
\end{formulabox}

\subsection{Divergence}

Divergence measures source-like or sink-like behaviour. A positive divergence at a point suggests local outward spreading. A negative divergence suggests local inward concentration. Zero divergence does not mean the vector field is zero; it means the field has no net local source or sink.

\subsection{Curl}

Curl measures local rotational tendency. A non-zero curl indicates that a small paddle wheel placed in the field would tend to rotate. A zero curl means the field is irrotational locally, although this does not automatically settle global questions in domains with holes.

\subsection{Directional Derivatives}

For a scalar field \(f\), the directional derivative in a unit direction \(\hat{\mathbf{u}}\) is
\[
    D_{\hat{\mathbf{u}}}f=\nabla f\cdot\hat{\mathbf{u}}.
\]
For a vector field, the derivative along \(\hat{\mathbf{u}}\) is
\[
    (\hat{\mathbf{u}}\cdot\nabla)\mathbf{A}.
\]
This gives another vector, not a scalar.

\begin{formulaboxnosplit}[Directional Derivative Formulae]
    If
    \[
        \hat{\mathbf{u}}=u_x\ihat+u_y\jhat+u_z\khat,\qquad \normval{\hat{\mathbf{u}}}=1,
    \]
    then
    \[
        D_{\hat{\mathbf{u}}}f=u_x\frac{\partial f}{\partial x}+u_y\frac{\partial f}{\partial y}+u_z\frac{\partial f}{\partial z}.
    \]

    \[
        \max_{\normval{\hat{\mathbf{u}}}=1}D_{\hat{\mathbf{u}}}f=\normval{\nabla f},
        \qquad
        \text{direction of maximum increase}=\frac{\nabla f}{\normval{\nabla f}}.
    \]

    \[
        (\hat{\mathbf{u}}\cdot\nabla)\mathbf{A}=
        \left(u_x\partial_x+u_y\partial_y+u_z\partial_z\right)\mathbf{A}.
    \]
\end{formulaboxnosplit}

\subsection{Product and Composition Identities}

These identities are worth learning in families. Notice how scalar multiplication, dot products and cross products change the type of the output.

\begin{formulabox}[Core Vector Identities]
    \[
        \nabla(fg)=f\nabla g+g\nabla f.
    \]

    \[
        \nabla\cdot(f\mathbf{A})=\nabla f\cdot\mathbf{A}+f\nabla\cdot\mathbf{A}.
    \]

    \[
        \nabla\times(f\mathbf{A})=\nabla f\times\mathbf{A}+f\nabla\times\mathbf{A}.
    \]

    \[
        \nabla(\mathbf{A}\cdot\mathbf{B})=(\mathbf{A}\cdot\nabla)\mathbf{B}+(\mathbf{B}\cdot\nabla)\mathbf{A}+\mathbf{A}\times(\nabla\times\mathbf{B})+\mathbf{B}\times(\nabla\times\mathbf{A}).
    \]

    \[
        \nabla\cdot(\mathbf{A}\times\mathbf{B})=\mathbf{B}\cdot(\nabla\times\mathbf{A})-\mathbf{A}\cdot(\nabla\times\mathbf{B}).
    \]

    \[
        \nabla\times(\mathbf{A}\times\mathbf{B})=\mathbf{A}(\nabla\cdot\mathbf{B})-\mathbf{B}(\nabla\cdot\mathbf{A})+(\mathbf{B}\cdot\nabla)\mathbf{A}-(\mathbf{A}\cdot\nabla)\mathbf{B}.
    \]

    \[
        \nabla\times(\nabla f)=\mathbf{0},\qquad \nabla\cdot(\nabla\times\mathbf{A})=0.
    \]

    \[
        \nabla\times(\nabla\times\mathbf{A})=\nabla(\nabla\cdot\mathbf{A})-\nabla^2\mathbf{A}.
    \]
\end{formulabox}

\begin{warningboxnosplit}[Common Mistakes]
    \begin{enumerate}
        \item Do not say that \(\nabla\mathbf{A}\) is automatically a vector. It is usually a second-order tensor.
        \item Do not confuse \(\nabla\cdot\mathbf{A}\) with \(\nabla\mathbf{A}\). Divergence is only the trace of the vector gradient in Cartesian coordinates.
        \item Do not use the Cartesian curl formula unchanged in cylindrical or spherical coordinates.
        \item Do not forget to normalise the direction vector before computing a directional derivative.
        \item Do not treat \(\nabla\) as an ordinary vector in every context. It is a differential operator.
    \end{enumerate}
\end{warningboxnosplit}

\begin{practicebox}[Quick Drills and Practice: Gradient, Divergence and Curl]
    \textbf{Level A:\,Formula recall}
    \begin{enumerate}
        \item Complete: \(\nabla=\ihat\quad\blank[1.6cm]+\jhat\quad\blank[1.6cm]+\khat\quad\blank[1.6cm]\).
        \item Complete: \(\nabla f=\blank[3cm]\).
        \item Complete: \(\nabla\cdot\mathbf{A}=\blank[3cm]\).
        \item Complete: \({(\nabla\times\mathbf{A})}_x=\blank[3cm]\).
        \item Complete: \({(\nabla\times\mathbf{A})}_y=\blank[3cm]\).
        \item Complete: \({(\nabla\times\mathbf{A})}_z=\blank[3cm]\).
        \item Complete: \(D_{\hat{\mathbf{u}}}f=\blank[3cm]\).
        \item Complete: \(\nabla\times(\nabla f)=\blank[2cm]\).
        \item Complete: \(\nabla\cdot(\nabla\times\mathbf{A})=\blank[2cm]\).
        \item Complete: \(\nabla\times(\nabla\times\mathbf{A})=\blank[6cm]\).
    \end{enumerate}

    \textbf{Level B:\,Direct computation}
    \begin{enumerate}
        \item Compute \(\nabla f\) for \(f=x^2+3y^2-z^2\).
        \item Compute \(\nabla f\) for \(f=xy^2+z\sin x\).
        \item Compute \(\nabla\cdot\mathbf{A}\) for \(\mathbf{A}=x\ihat+y\jhat+z\khat\).
        \item Compute \(\nabla\cdot\mathbf{A}\) for \(\mathbf{A}=yz\ihat+xz\jhat+xy\khat\).
        \item Compute \(\nabla\times\mathbf{A}\) for \(\mathbf{A}=(-y)\ihat+x\jhat+0\khat\).
        \item Compute \(\nabla\times\mathbf{A}\) for \(\mathbf{A}=z\ihat+x\jhat+y\khat\).
        \item For \(f=x^2y+yz^2\), find \(D_{\hat{\mathbf{u}}}f\) at \((1,2,1)\) in the direction \(\mathbf{u}=2\ihat-\jhat+2\khat\).
        \item For \(\mathbf{A}=x^2\ihat+xy\jhat+z^2\khat\), compute \(\nabla\mathbf{A}\), \(\nabla\cdot\mathbf{A}\), and \(\nabla\times\mathbf{A}\).
    \end{enumerate}

    \textbf{Level C:\,Interpretation and multi-step questions}
    \begin{enumerate}
        \item A vector field has \(\nabla\cdot\mathbf{A}>0\) near a point. Explain the local behaviour in words.
        \item Give a vector field with zero divergence but non-zero curl.
        \item Give a vector field with non-zero divergence but zero curl.
        \item If \(\nabla f=\mathbf{0}\) throughout a connected region, what can be said about \(f\)?
        \item Find the direction of maximum increase of \(f=x^2+y^2+z^2\) at \((1,1,2)\).
        \item Find the maximum directional derivative of \(f=xy+yz+zx\) at \((1,2,3)\).
    \end{enumerate}

    \textbf{Level D:\,Proofs and identities}
    \begin{enumerate}
        \item Prove in Cartesian components that \(\nabla\times(\nabla f)=\mathbf{0}\), assuming mixed partial derivatives are equal.
        \item Prove in Cartesian components that \(\nabla\cdot(\nabla\times\mathbf{A})=0\).
        \item Verify the curl-curl identity for \(\mathbf{A}=x^2\ihat+y^2\jhat+z^2\khat\).
        \item Derive the product rule for \(\nabla\cdot(f\mathbf{A})\).
        \item Derive the product rule for \(\nabla\times(f\mathbf{A})\).
    \end{enumerate}

    \textbf{Level E:\,Challenge applications}
    \begin{enumerate}
        \item In fluid mechanics, \(\mathbf{v}\) is a velocity field. Interpret \(\nabla\cdot\mathbf{v}=0\).
        \item In electromagnetism, \(\nabla\times\mathbf{E}=\mathbf{0}\) in electrostatics. What does this imply about a scalar potential?
        \item For \(\mathbf{v}=(-\omega y)\ihat+(\omega x)\jhat\), compute divergence and curl, then interpret the motion.
        \item For \(\mathbf{v}=a x\ihat+a y\jhat+a z\khat\), compute divergence and curl, then interpret the motion.
    \end{enumerate}
\end{practicebox}

\ifdefined\OTCOMBINED%
\else
\end{document}
\fi

\ifdefined\OTCOMBINED%
\else
\documentclass{otscience}
% ============================================================
% Shared setup for OT Science modular workbook parts
% Place these files in the same folder as your fixed otscience.cls
% ============================================================

\makeatletter
\makeatother
\usepackage{multicol}
\usepackage{longtable}
\usepackage{makecell}
\usepackage{pdflscape}
\usepackage{bm}
\providecommand{\blank}[1][3cm]{\rule{#1}{0.4pt}}
\providecommand{\bigblank}{\rule{7cm}{0.4pt}}
\providecommand{\componentblank}{\blank[2.5cm]}
\providecommand{\R}{\mathbb{R}}
\providecommand{\C}{\mathbb{C}}
\providecommand{\ii}{\mathrm{i}}
\providecommand{\ee}{\mathrm{e}}
\providecommand{\dd}{\mathrm{d}}
\providecommand{\grad}{\nabla}
\providecommand{\divergence}{\nabla\!\cdot}
\providecommand{\curl}{\nabla\!\times}
\providecommand{\lap}{\nabla^2}
\providecommand{\ihat}{\hat{\mathbf{i}}}
\providecommand{\jhat}{\hat{\mathbf{j}}}
\providecommand{\khat}{\hat{\mathbf{k}}}
\providecommand{\er}{\hat{\mathbf{e}}_r}
\providecommand{\etheta}{\hat{\mathbf{e}}_\theta}
\providecommand{\ephi}{\hat{\mathbf{e}}_\phi}
\providecommand{\erho}{\hat{\boldsymbol{\rho}}}
\providecommand{\vA}{\mathbf{A}}
\providecommand{\vB}{\mathbf{B}}
\providecommand{\va}{\mathbf{a}}
\providecommand{\vb}{\mathbf{b}}
\providecommand{\vc}{\mathbf{c}}
\providecommand{\vx}{\mathbf{x}}
\providecommand{\vr}{\mathbf{r}}
\providecommand{\matA}{\mathbf{A}}
\providecommand{\matB}{\mathbf{B}}
\providecommand{\tr}{\operatorname{tr}}
\providecommand{\cof}{\operatorname{cof}}
\providecommand{\dev}{\operatorname{dev}}
\providecommand{\diag}{\operatorname{diag}}
\providecommand{\questionline}[1][5cm]{\rule{#1}{0.4pt}}
\setlength{\parindent}{0pt}
\setlength{\parskip}{4pt}
\renewcommand{\arraystretch}{1.25}

% Fallbacks in case the current otscience.cls has not defined nosplit boxes.
\makeatletter
\@ifundefined{formulaboxnosplit}{\newenvironment{formulaboxnosplit}[1][Formula]{\begin{formulabox}[#1]}{\end{formulabox}}}{}
\@ifundefined{definitionboxnosplit}{\newenvironment{definitionboxnosplit}[1][Definition]{\begin{definitionbox}[#1]}{\end{definitionbox}}}{}
\@ifundefined{summaryboxnosplit}{\newenvironment{summaryboxnosplit}[1][Summary]{\begin{summarybox}[#1]}{\end{summarybox}}}{}
\@ifundefined{warningboxnosplit}{\newenvironment{warningboxnosplit}[1][Warning]{\begin{warningbox}[#1]}{\end{warningbox}}}{}
\@ifundefined{exampleboxnosplit}{\newenvironment{exampleboxnosplit}[1][Example]{\begin{examplebox}[#1]}{\end{examplebox}}}{}
\makeatother

\newenvironment{practicebox}[1][Quick Drills and Practice]{\begin{examplebox}[#1]}{\end{examplebox}}
\newenvironment{practiceboxnosplit}[1][Quick Drills and Practice]{\begin{exampleboxnosplit}[#1]}{\end{exampleboxnosplit}}

\newcommand{\standalonetitle}[3]{%
    \title{\textbf{#1}\\\large #2}
    \author{OT Science Notes}
    \date{#3}
}

\standalonetitle{Coordinate Systems and Tensor Operations}{Part 2 of the OT Science Formula Completion Workbook}{\DTMtoday}
\begin{document}
\maketitle
\tableofcontents
\newpage
\fi

% 02_coordinate_systems_tensor_operations.tex
\section{Coordinate Systems and Tensor Operations}
\sectionmark{Coordinate Systems}

Coordinate systems are not just different labels for the same point. They also change how basis vectors, scale factors, area elements, volume elements and differential operators are written. Tensor operations help us keep the geometry clear when the coordinates become less simple than rectangular coordinates.

\subsection{Rectangular, Cylindrical and Spherical Coordinates}

Rectangular coordinates are the easiest because the basis vectors are fixed. Cylindrical and spherical coordinates are more natural for circular, axial, radial and angular symmetry, but their basis vectors change with position.

\begin{formulabox}[Coordinate Transformations]
    \[
        \text{Cartesian:}\qquad (x,y,z).
    \]

    \[
        \text{Cylindrical:}\qquad x=\rho\cos\phi,\quad y=\rho\sin\phi,\quad z=z.
    \]

    \[
        \rho=\sqrt{x^2+y^2},\qquad \phi=\operatorname{atan2}(y,x),\qquad z=z.
    \]

    \[
        \text{Spherical:}\qquad x=r\sin\theta\cos\phi,
        \quad y=r\sin\theta\sin\phi,
        \quad z=r\cos\theta.
    \]

    \[
        r=\sqrt{x^2+y^2+z^2},\qquad
        \theta=\cos^{-1}\left(\frac{z}{r}\right),\qquad
        \phi=\operatorname{atan2}(y,x).
    \]

    \[
        \rho=r\sin\theta,
        \qquad z=r\cos\theta,
        \qquad r=\sqrt{\rho^2+z^2}.
    \]
\end{formulabox}

\subsection{Basis Vectors and Scale Factors}

A coordinate basis is obtained by differentiating the position vector with respect to the coordinates. In orthogonal coordinates, the scale factors tell us how much physical length corresponds to a small coordinate change.

\begin{formulabox}[Basis Vectors, Scale Factors and Metric Matrices]
    For orthogonal coordinates \((q^1,q^2,q^3)\),
    \[
        \mathbf{g}_i=\frac{\partial\mathbf{r}}{\partial q^i},
        \qquad
        h_i=\normval{\mathbf{g}_i},
        \qquad
        \hat{\mathbf{e}}_i=\frac{\mathbf{g}_i}{h_i}.
    \]

    \[
        ds^2=h_1^2{(dq^1)}^2+h_2^2{(dq^2)}^2+h_3^2{(dq^3)}^2.
    \]

    \[
        [g_{ij}]=\begin{pmatrix}
            h_1^2 & 0     & 0     \\
            0     & h_2^2 & 0     \\
            0     & 0     & h_3^2
        \end{pmatrix},
        \qquad
        [g^{ij}]=\begin{pmatrix}
            1/h_1^2 & 0       & 0       \\
            0       & 1/h_2^2 & 0       \\
            0       & 0       & 1/h_3^2
        \end{pmatrix}.
    \]

    \[
        \text{Cartesian:}\qquad h_x=h_y=h_z=1,
        \qquad ds^2=dx^2+dy^2+dz^2.
    \]

    \[
        \text{Cylindrical:}\qquad h_\rho=1,
        \qquad h_\phi=\rho,
        \qquad h_z=1,
        \qquad ds^2=d\rho^2+\rho^2d\phi^2+dz^2.
    \]

    \[
        \text{Spherical:}\qquad h_r=1,
        \qquad h_\theta=r,
        \qquad h_\phi=r\sin\theta,
    \]
    \[
        ds^2=dr^2+r^2d\theta^2+r^2\sin^2\theta\,d\phi^2.
    \]
\end{formulabox}

\subsection{Volume and Area Elements}

The volume element is controlled by the product of the scale factors. This is why cylindrical and spherical integrals contain the extra factors \(\rho\) and \(r^2\sin\theta\).

\begin{formulaboxnosplit}[Differential Elements]
    \[
        dV=h_1h_2h_3\,dq^1dq^2dq^3.
    \]

    \[
        dV_{\text{cart}}=dx\,dy\,dz.
    \]

    \[
        dV_{\text{cyl}}=\rho\,d\rho\,d\phi\,dz.
    \]

    \[
        dV_{\text{sph}}=r^2\sin\theta\,dr\,d\theta\,d\phi.
    \]

    For orthogonal coordinates, the coordinate surface elements are
    \[
        dS_1=h_2h_3\,dq^2dq^3,
        \qquad
        dS_2=h_3h_1\,dq^3dq^1,
        \qquad
        dS_3=h_1h_2\,dq^1dq^2.
    \]
\end{formulaboxnosplit}

\subsection{Coordinate Basis, Physical Components and Tensor Components}

A vector can be represented in coordinate basis vectors or in unit basis vectors. In many engineering texts, cylindrical and spherical vector fields are written using unit basis vectors, so the listed components are physical components. In tensor notation, however, components may be covariant or contravariant depending on the basis used.

\begin{formulabox}[Basic Tensor Operations]
    For a second-order tensor \(A_{ij}\),
    \[
        {(\mathbf{A}\mathbf{x})}_i=A_{ij}x_j,
        \qquad
        \tr(\mathbf{A})=A_{ii},
        \qquad
        \mathbf{A}:\mathbf{B}=A_{ij}B_{ij}.
    \]

    \[
        {(\mathbf{A}\mathbf{B})}_{ij}=A_{ik}B_{kj}.
    \]

    \[
        \mathbf{x}\cdot\mathbf{A}\mathbf{x}=x_{i}A_{ij}x_j.
    \]

    \[
        \nabla\mathbf{A}\quad\text{for a vector field gives}\quad {(\nabla\mathbf{A})}_{ij}=\partial_j A_i.
    \]

    \[
        \nabla\cdot\mathbf{T}\quad\text{for a second-order tensor can be written}\quad {(\nabla\cdot\mathbf{T})}_i=\partial_{j}T_{ij}.
    \]
\end{formulabox}

\subsection{Gradient, Divergence and Curl in Orthogonal Coordinates}

The following formulas use physical components \(A_1,A_2,A_3\) along unit vectors \(\hat{\mathbf{e}}_1,\hat{\mathbf{e}}_2,\hat{\mathbf{e}}_3\).

\begingroup
\small
\setlength{\arraycolsep}{3pt}
\[
    \nabla\times\mathbf{A}
    =
    \frac{1}{h_1h_2h_3}
    \begin{vmatrix}
        h_1\hat{\mathbf{e}}_1 & h_2\hat{\mathbf{e}}_2 & h_3\hat{\mathbf{e}}_3 \\
        \partial_{q^1}        & \partial_{q^2}        & \partial_{q^3}        \\
        h_1A_1                & h_2A_2                & h_3A_3
    \end{vmatrix}.
\]
\endgroup

\subsection{Cylindrical Differential Operators}

\begin{formulabox}[Cylindrical: \((\rho,\phi,z)\)]
    \[
        \nabla f=\erho\frac{\partial f}{\partial\rho}+\hat{\boldsymbol{\phi}}\frac{1}{\rho}\frac{\partial f}{\partial\phi}+\khat\frac{\partial f}{\partial z}.
    \]

    \[
        \nabla\cdot\mathbf{A}=\frac{1}{\rho}\frac{\partial}{\partial\rho}(\rho A_\rho)+\frac{1}{\rho}\frac{\partial A_\phi}{\partial\phi}+\frac{\partial A_z}{\partial z}.
    \]

    \[
        \nabla\times\mathbf{A}=
        \erho\left(\frac{1}{\rho}\frac{\partial A_z}{\partial\phi}-\frac{\partial A_\phi}{\partial z}\right)
        +\hat{\boldsymbol{\phi}}\left(\frac{\partial A_\rho}{\partial z}-\frac{\partial A_z}{\partial\rho}\right)
        +\khat\left(\frac{1}{\rho}\frac{\partial}{\partial\rho}(\rho A_\phi)-\frac{1}{\rho}\frac{\partial A_\rho}{\partial\phi}\right).
    \]
\end{formulabox}

\subsection{Spherical Differential Operators}

\begin{formulabox}[Spherical: \((r,\theta,\phi)\)]
\[
\begin{aligned}
    \nabla\cdot\mathbf{A}
    ={}&
    \frac{1}{r^2}
    \frac{\partial}{\partial r}(r^2A_r)
    \\[2pt]
    &+
    \frac{1}{r\sin\theta}
    \frac{\partial}{\partial\theta}
    (\sin\theta A_\theta)
    \\[2pt]
    &+
    \frac{1}{r\sin\theta}
    \frac{\partial A_\phi}{\partial\phi}.
\end{aligned}
\]

\[
\begin{aligned}
    \nabla\times\mathbf{A}
    ={}&
    \er\frac{1}{r\sin\theta}
    \left[
        \frac{\partial}{\partial\theta}(\sin\theta A_\phi)
        -
        \frac{\partial A_\theta}{\partial\phi}
    \right]
    \\[3pt]
    &+
    \etheta\frac{1}{r}
    \left[
        \frac{1}{\sin\theta}
        \frac{\partial A_r}{\partial\phi}
        -
        \frac{\partial}{\partial r}(rA_\phi)
    \right]
    \\[3pt]
    &+
    \ephi\frac{1}{r}
    \left[
        \frac{\partial}{\partial r}(rA_\theta)
        -
        \frac{\partial A_r}{\partial\theta}
    \right].
\end{aligned}
\]
\end{formulabox}

\begin{practicebox}[Quick Drills and Practice: Coordinates and Tensor Operations]
    \textbf{Level A:\,Coordinate conversion}
    \begin{multicols}{2}
        \begin{enumerate}
            \item Convert \((x,y,z)=(3,4,5)\) to cylindrical coordinates.
            \item Convert \((x,y,z)=(1,1,\sqrt2)\) to spherical coordinates.
            \item Convert \((\rho,\phi,z)=(2,\pi/3,5)\) to Cartesian coordinates.
            \item Convert \((r,\theta,\phi)=(4,\pi/2,\pi/6)\) to Cartesian coordinates.
            \item Fill in: \(\rho=\blank[2cm]\), \(z=\blank[2cm]\), \(r=\blank[2cm]\).
            \item Fill in: \(dV_{\text{cyl}}=\blank[3cm]\).
            \item Fill in: \(dV_{\text{sph}}=\blank[3cm]\).
            \item Fill in the spherical scale factors.
        \end{enumerate}
    \end{multicols}

    \textbf{Level B:\,Scale factors and metric matrices}
    \begin{enumerate}
        \item Starting from \(\mathbf{r}=\rho\cos\phi\,\ihat+\rho\sin\phi\,\jhat+z\khat\), derive \(h_\rho,h_\phi,h_z\).
        \item Starting from \(\mathbf{r}=r\sin\theta\cos\phi\,\ihat+r\sin\theta\sin\phi\,\jhat+r\cos\theta\,\khat\), derive \(h_r,h_\theta,h_\phi\).
        \item Write the metric matrix \([g_{ij}]\) for cylindrical coordinates.
        \item Write the metric matrix \([g_{ij}]\) for spherical coordinates.
        \item Find \(ds^2\) for cylindrical coordinates and explain why the term \(\rho^2d\phi^2\) appears.
        \item Find \(ds^2\) for spherical coordinates and explain the term \(r^2\sin^2\theta\,d\phi^2\).
    \end{enumerate}

    \textbf{Level C:\,Differential operators}
    \begin{enumerate}
        \item Compute \(\nabla f\) in cylindrical coordinates for \(f=\rho^2z\sin\phi\).
        \item Compute \(\nabla f\) in spherical coordinates for \(f=r^2\cos\theta\).
        \item Compute \(\nabla\cdot\mathbf{A}\) in cylindrical coordinates for \(\mathbf{A}=\rho^2\erho+\rho z\hat{\boldsymbol{\phi}}+z^2\khat\).
        \item Compute \(\nabla\cdot\mathbf{A}\) in spherical coordinates for \(\mathbf{A}=r^2\er+r\sin\theta\,\etheta+0\ephi\).
        \item Compute \(\nabla\times\mathbf{A}\) in cylindrical coordinates for \(\mathbf{A}=0\erho+\rho^2\hat{\boldsymbol{\phi}}+0\khat\).
        \item Compute \(\nabla\times\mathbf{A}\) in spherical coordinates for \(\mathbf{A}=0\er+0\etheta+r\sin\theta\,\ephi\).
    \end{enumerate}

    \textbf{Level D:\,Tensor operations}
    \begin{enumerate}
        \item Given \(A_{ij}=\begin{pmatrix}1&2&0\\0&3&1\\4&0&2\end{pmatrix}\), compute \(A_{ij}x_j\) for \(\mathbf{x}=(1,2,3)\).
        \item Compute \(\tr(A)\), \(A:A\), and \(A^2\) for the matrix above.
        \item If \(T_{ij}=x_{i}x_j\), compute \(\partial_{j}T_{ij}\) in three-dimensional Cartesian coordinates.
        \item For \(\mathbf{A}=x^2\ihat+xy\jhat+yz\khat\), compute \({(\nabla\mathbf{A})}_{ij}\).
        \item Show explicitly that \(\nabla\cdot\mathbf{A}=\tr(\nabla\mathbf{A})\) for the same vector field.
    \end{enumerate}
\end{practicebox}

\ifdefined\OTCOMBINED%
\else
\end{document}
\fi

\ifdefined\OTCOMBINED%
\else
\documentclass{otscience}
% ============================================================
% Shared setup for OT Science modular workbook parts
% Place these files in the same folder as your fixed otscience.cls
% ============================================================

\makeatletter
\makeatother
\usepackage{multicol}
\usepackage{longtable}
\usepackage{makecell}
\usepackage{pdflscape}
\usepackage{bm}
\providecommand{\blank}[1][3cm]{\rule{#1}{0.4pt}}
\providecommand{\bigblank}{\rule{7cm}{0.4pt}}
\providecommand{\componentblank}{\blank[2.5cm]}
\providecommand{\R}{\mathbb{R}}
\providecommand{\C}{\mathbb{C}}
\providecommand{\ii}{\mathrm{i}}
\providecommand{\ee}{\mathrm{e}}
\providecommand{\dd}{\mathrm{d}}
\providecommand{\grad}{\nabla}
\providecommand{\divergence}{\nabla\!\cdot}
\providecommand{\curl}{\nabla\!\times}
\providecommand{\lap}{\nabla^2}
\providecommand{\ihat}{\hat{\mathbf{i}}}
\providecommand{\jhat}{\hat{\mathbf{j}}}
\providecommand{\khat}{\hat{\mathbf{k}}}
\providecommand{\er}{\hat{\mathbf{e}}_r}
\providecommand{\etheta}{\hat{\mathbf{e}}_\theta}
\providecommand{\ephi}{\hat{\mathbf{e}}_\phi}
\providecommand{\erho}{\hat{\boldsymbol{\rho}}}
\providecommand{\vA}{\mathbf{A}}
\providecommand{\vB}{\mathbf{B}}
\providecommand{\va}{\mathbf{a}}
\providecommand{\vb}{\mathbf{b}}
\providecommand{\vc}{\mathbf{c}}
\providecommand{\vx}{\mathbf{x}}
\providecommand{\vr}{\mathbf{r}}
\providecommand{\matA}{\mathbf{A}}
\providecommand{\matB}{\mathbf{B}}
\providecommand{\tr}{\operatorname{tr}}
\providecommand{\cof}{\operatorname{cof}}
\providecommand{\dev}{\operatorname{dev}}
\providecommand{\diag}{\operatorname{diag}}
\providecommand{\questionline}[1][5cm]{\rule{#1}{0.4pt}}
\setlength{\parindent}{0pt}
\setlength{\parskip}{4pt}
\renewcommand{\arraystretch}{1.25}

% Fallbacks in case the current otscience.cls has not defined nosplit boxes.
\makeatletter
\@ifundefined{formulaboxnosplit}{\newenvironment{formulaboxnosplit}[1][Formula]{\begin{formulabox}[#1]}{\end{formulabox}}}{}
\@ifundefined{definitionboxnosplit}{\newenvironment{definitionboxnosplit}[1][Definition]{\begin{definitionbox}[#1]}{\end{definitionbox}}}{}
\@ifundefined{summaryboxnosplit}{\newenvironment{summaryboxnosplit}[1][Summary]{\begin{summarybox}[#1]}{\end{summarybox}}}{}
\@ifundefined{warningboxnosplit}{\newenvironment{warningboxnosplit}[1][Warning]{\begin{warningbox}[#1]}{\end{warningbox}}}{}
\@ifundefined{exampleboxnosplit}{\newenvironment{exampleboxnosplit}[1][Example]{\begin{examplebox}[#1]}{\end{examplebox}}}{}
\makeatother

\newenvironment{practicebox}[1][Quick Drills and Practice]{\begin{examplebox}[#1]}{\end{examplebox}}
\newenvironment{practiceboxnosplit}[1][Quick Drills and Practice]{\begin{exampleboxnosplit}[#1]}{\end{exampleboxnosplit}}

\newcommand{\standalonetitle}[3]{%
    \title{\textbf{#1}\\\large #2}
    \author{OT Science Notes}
    \date{#3}
}

\standalonetitle{Index Notation, Kronecker Delta and Levi-Civita Symbol}{Part 3 of the OT Science Formula Completion Workbook}{\DTMtoday}
\begin{document}
\maketitle
\tableofcontents
\newpage
\fi

% 03_index_delta_levicivita.tex
\section{Index Notation, Kronecker Delta and Levi-Civita Symbol}
\sectionmark{Index Notation}

Index notation compresses vector and tensor formulas into short expressions. It is especially useful when proving vector identities. The price is that every index must be handled carefully.

\subsection{Einstein Summation Convention}

When an index is repeated once as a subscript or repeated in a product in Euclidean notation, it is summed over its full range. In three-dimensional Cartesian coordinates, the range is usually \(1,2,3\).

\begin{formulaboxnosplit}[Einstein Summation]
    \[
        a_i b_i=a_1b_1+a_2b_2+a_3b_3.
    \]

    \[
        A_{ij}b_j=A_{i1}b_1+A_{i2}b_2+A_{i3}b_3.
    \]

    \[
        A_{ij}B_{ij}=\sum_{i=1}^{3}\sum_{j=1}^{3}A_{ij}B_{ij}.
    \]

    \[
        {(A B)}_{ij}=A_{ik}B_{kj}.
    \]
\end{formulaboxnosplit}

A free index appears on both sides of an equation and is not summed. A dummy index is repeated and summed. Dummy indices can be renamed as long as no conflict is created.

\subsection{Kronecker Delta}

The Kronecker delta is the identity tensor in component form. It replaces one index by another.

\begin{formulabox}[Kronecker Delta Identities]
    \[
        \delta_{ij}=\begin{cases}
            1, & i=j,     \\
            0, & i\neq j.
        \end{cases}
    \]

    \[
        \delta_{ij}a_j=a_i,
        \qquad
        \delta_{ij}a_i=a_j.
    \]

    \[
        \delta_{ij}\delta_{jk}=\delta_{ik},
        \qquad
        \delta_{ii}=3.
    \]

    \[
        \delta_{ij}A_{ij}=A_{ii}=\tr(A).
    \]

    \[
        \partial_i x_j=\frac{\partial x_j}{\partial x_i}=\delta_{ij}.
    \]

    \[
        \partial_i x_i=3,
        \qquad
        \partial_i r=\frac{x_i}{r},
        \qquad
        \partial_i(r^2)=2x_i.
    \]
\end{formulabox}

\subsection{Levi-Civita Symbol}

The Levi-Civita symbol tracks orientation and cross products. It is zero whenever any two indices are repeated.

\begin{formulabox}[Levi-Civita Identities]
    \[
        \varepsilon_{ijk}=\begin{cases}
            +1, & (i,j,k)\text{ is an even permutation of }(1,2,3), \\
            -1, & (i,j,k)\text{ is an odd permutation of }(1,2,3),  \\
            0,  & \text{two indices are repeated}.
        \end{cases}
    \]

    \[
        {(\mathbf{a}\times\mathbf{b})}_i=\varepsilon_{ijk}a_{j}b_k.
    \]

    \[
        \varepsilon_{ijk}\varepsilon_{ilm}=\delta_{jl}\delta_{km}-\delta_{jm}\delta_{kl}.
    \]

    \[
        \varepsilon_{ijk}\varepsilon_{ijm}=2\delta_{km},
        \qquad
        \varepsilon_{ijk}\varepsilon_{ijk}=6.
    \]

    \[
        \varepsilon_{ijk}\varepsilon_{lmn}=\begin{vmatrix}
            \delta_{il} & \delta_{im} & \delta_{in} \\
            \delta_{jl} & \delta_{jm} & \delta_{jn} \\
            \delta_{kl} & \delta_{km} & \delta_{kn}
        \end{vmatrix}.
    \]
\end{formulabox}

\subsection{Vector Calculus in Index Notation}

\begin{formulabox}[Vector Operations in Index Form]
    \[
        {(\nabla f)}_i=\partial_i f.
    \]

    \[
        \nabla\cdot\mathbf{A}=\partial_{i}A_i.
    \]

    \[
        {(\nabla\times\mathbf{A})}_i=\varepsilon_{ijk}\partial_{j}A_k.
    \]

    \[
        \nabla^2 f=\partial_i\partial_i f.
    \]

    \[
        {(\nabla\mathbf{A})}_{ij}=\partial_{j}A_i.
    \]

    \[
        {(\nabla\cdot\mathbf{T})}_i=\partial_{j}T_{ij}.
    \]
\end{formulabox}

\subsection{Proof Pattern: Curl of Curl}

The index method is powerful because the proof becomes a controlled substitution. Starting with
\[
    {(\nabla\times(\nabla\times\mathbf{A}))}_i=\varepsilon_{ijk}\partial_j{(\nabla\times\mathbf{A})}_k,
\]
use
\[
    {(\nabla\times\mathbf{A})}_k=\varepsilon_{klm}\partial_{l}A_m.
\]
Then
\[
    {(\nabla\times(\nabla\times\mathbf{A}))}_i=\varepsilon_{ijk}\varepsilon_{klm}\partial_j\partial_{l}A_m.
\]
Using the contraction identity gives the familiar result
\[
    \nabla\times(\nabla\times\mathbf{A})=\nabla(\nabla\cdot\mathbf{A})-\nabla^2\mathbf{A}.
\]

\begin{warningboxnosplit}[Index-Notation Mistakes]
    \begin{enumerate}
        \item Do not repeat the same index more than twice in one term unless you have explicitly defined what it means.
        \item Do not change a free index on one side of an equation without changing it on the other side.
        \item Dummy indices may be renamed, but free indices may not.
        \item The identity \(\varepsilon_{ijk}\varepsilon_{ilm}\) depends on which index is contracted. Check the repeated index carefully.
    \end{enumerate}
\end{warningboxnosplit}

\begin{practicebox}[Quick Drills and Practice: Index Notation]
    \textbf{Level A:\,Expansion and simplification}
    \begin{multicols}{2}
        \begin{enumerate}
            \item Expand \(a_{i}b_i\) in three dimensions.
            \item Expand \(A_{ij}b_j\) for \(i=1\).
            \item Simplify \(\delta_{ij}a_j\).
            \item Simplify \(\delta_{ij}A_{jk}\).
            \item Simplify \(\delta_{ij}\delta_{jk}\delta_{kl}\).
            \item Simplify \(\delta_{ij}A_{ij}\).
            \item Compute \(\delta_{ii}\) in three dimensions.
            \item Compute \(\partial_i x_j\).
            \item Compute \(\partial_i x_i\).
            \item Compute \(\partial_i(x_{i}x_j)\).
        \end{enumerate}
    \end{multicols}

    \textbf{Level B:\,Levi-Civita}
    \begin{enumerate}
        \item Evaluate \(\varepsilon_{123}\), \(\varepsilon_{132}\), \(\varepsilon_{221}\), \(\varepsilon_{312}\).
        \item Expand \({(\mathbf{a}\times\mathbf{b})}_1=\varepsilon_{1jk}a_{j}b_k\).
        \item Show that \(\mathbf{a}\times\mathbf{b}=-\mathbf{b}\times\mathbf{a}\) using \(\varepsilon_{ijk}\).
        \item Simplify \(\varepsilon_{ijk}\varepsilon_{ijm}\).
        \item Simplify \(\varepsilon_{ijk}\varepsilon_{ijk}\).
        \item Use \(\varepsilon_{ijk}\varepsilon_{ilm}\) to simplify \((\mathbf{a}\times\mathbf{b})\cdot(\mathbf{c}\times\mathbf{d})\).
    \end{enumerate}

    \textbf{Level C:\,Vector identities}
    \begin{enumerate}
        \item Prove \(\nabla\cdot(\nabla\times\mathbf{A})=0\) using index notation.
        \item Prove \(\nabla\times(\nabla f)=\mathbf{0}\) using index notation.
        \item Prove \(\nabla\times(\nabla\times\mathbf{A})=\nabla(\nabla\cdot\mathbf{A})-\nabla^2\mathbf{A}\).
        \item Prove \(\mathbf{a}\times(\mathbf{b}\times\mathbf{c})=\mathbf{b}(\mathbf{a}\cdot\mathbf{c})-\mathbf{c}(\mathbf{a}\cdot\mathbf{b})\).
        \item Prove \((\mathbf{a}\times\mathbf{b})\cdot(\mathbf{c}\times\mathbf{d})=(\mathbf{a}\cdot\mathbf{c})(\mathbf{b}\cdot\mathbf{d})-(\mathbf{a}\cdot\mathbf{d})(\mathbf{b}\cdot\mathbf{c})\).
    \end{enumerate}

    \textbf{Level D:\,Derivative identities}
    \begin{enumerate}
        \item Show that \(\partial_i r=x_i/r\), where \(r={(x_{j}x_j)}^{1/2}\).
        \item Show that \(\partial_i(1/r)=-x_i/r^3\).
        \item Compute \(\partial_i(x_i/r^3)\) for \(r\neq0\).
        \item Compute \(\partial_j(x_{i}x_j)\) in three dimensions.
        \item If \(T_{ij}=f\delta_{ij}\), compute \(\partial_{j}T_{ij}\).
    \end{enumerate}
\end{practicebox}

\ifdefined\OTCOMBINED%
\else
\end{document}
\fi

\ifdefined\OTCOMBINED%
\else
\documentclass{otscience}
% ============================================================
% Shared setup for OT Science modular workbook parts
% Place these files in the same folder as your fixed otscience.cls
% ============================================================

\makeatletter
\makeatother
\usepackage{multicol}
\usepackage{longtable}
\usepackage{makecell}
\usepackage{pdflscape}
\usepackage{bm}
\providecommand{\blank}[1][3cm]{\rule{#1}{0.4pt}}
\providecommand{\bigblank}{\rule{7cm}{0.4pt}}
\providecommand{\componentblank}{\blank[2.5cm]}
\providecommand{\R}{\mathbb{R}}
\providecommand{\C}{\mathbb{C}}
\providecommand{\ii}{\mathrm{i}}
\providecommand{\ee}{\mathrm{e}}
\providecommand{\dd}{\mathrm{d}}
\providecommand{\grad}{\nabla}
\providecommand{\divergence}{\nabla\!\cdot}
\providecommand{\curl}{\nabla\!\times}
\providecommand{\lap}{\nabla^2}
\providecommand{\ihat}{\hat{\mathbf{i}}}
\providecommand{\jhat}{\hat{\mathbf{j}}}
\providecommand{\khat}{\hat{\mathbf{k}}}
\providecommand{\er}{\hat{\mathbf{e}}_r}
\providecommand{\etheta}{\hat{\mathbf{e}}_\theta}
\providecommand{\ephi}{\hat{\mathbf{e}}_\phi}
\providecommand{\erho}{\hat{\boldsymbol{\rho}}}
\providecommand{\vA}{\mathbf{A}}
\providecommand{\vB}{\mathbf{B}}
\providecommand{\va}{\mathbf{a}}
\providecommand{\vb}{\mathbf{b}}
\providecommand{\vc}{\mathbf{c}}
\providecommand{\vx}{\mathbf{x}}
\providecommand{\vr}{\mathbf{r}}
\providecommand{\matA}{\mathbf{A}}
\providecommand{\matB}{\mathbf{B}}
\providecommand{\tr}{\operatorname{tr}}
\providecommand{\cof}{\operatorname{cof}}
\providecommand{\dev}{\operatorname{dev}}
\providecommand{\diag}{\operatorname{diag}}
\providecommand{\questionline}[1][5cm]{\rule{#1}{0.4pt}}
\setlength{\parindent}{0pt}
\setlength{\parskip}{4pt}
\renewcommand{\arraystretch}{1.25}

% Fallbacks in case the current otscience.cls has not defined nosplit boxes.
\makeatletter
\@ifundefined{formulaboxnosplit}{\newenvironment{formulaboxnosplit}[1][Formula]{\begin{formulabox}[#1]}{\end{formulabox}}}{}
\@ifundefined{definitionboxnosplit}{\newenvironment{definitionboxnosplit}[1][Definition]{\begin{definitionbox}[#1]}{\end{definitionbox}}}{}
\@ifundefined{summaryboxnosplit}{\newenvironment{summaryboxnosplit}[1][Summary]{\begin{summarybox}[#1]}{\end{summarybox}}}{}
\@ifundefined{warningboxnosplit}{\newenvironment{warningboxnosplit}[1][Warning]{\begin{warningbox}[#1]}{\end{warningbox}}}{}
\@ifundefined{exampleboxnosplit}{\newenvironment{exampleboxnosplit}[1][Example]{\begin{examplebox}[#1]}{\end{examplebox}}}{}
\makeatother

\newenvironment{practicebox}[1][Quick Drills and Practice]{\begin{examplebox}[#1]}{\end{examplebox}}
\newenvironment{practiceboxnosplit}[1][Quick Drills and Practice]{\begin{exampleboxnosplit}[#1]}{\end{exampleboxnosplit}}

\newcommand{\standalonetitle}[3]{%
    \title{\textbf{#1}\\\large #2}
    \author{OT Science Notes}
    \date{#3}
}

\standalonetitle{Covariant and Contravariant Bases, Metrics and General Curvilinear Calculus}{Part 4 of the OT Science Formula Completion Workbook}{\DTMtoday}
\begin{document}
\maketitle
\tableofcontents
\newpage
\fi

% 04_covariant_contravariant_metrics.tex
\section[
    Covariant and Contravariant Bases, Metrics and General Curvilinear Calculus
]{%
    Covariant and Contravariant Bases, Metrics\\
    and General Curvilinear Calculus%
}
\sectionmark{Metrics and Bases}

This part connects coordinate geometry with tensor notation. The key idea is that a vector itself is geometric, but its components depend on the basis used to describe it. Covariant and contravariant components are two related ways of recording the same geometric object.

\subsection{Coordinate Basis and Reciprocal Basis}

Let
\[
    \mathbf{r}=\mathbf{r}(q^1,q^2,q^3).
\]
The covariant basis vectors are tangent to coordinate curves:
\[
    \mathbf{e}_i=\frac{\partial\mathbf{r}}{\partial q^i}.
\]
The contravariant basis vectors form the reciprocal basis:
\[
    \mathbf{e}^i\cdot\mathbf{e}_j=\delta^i_j.
\]

\begin{formulaboxnosplit}[Covariant and Contravariant Bases]
    \[
        \mathbf{e}_i=\frac{\partial\mathbf{r}}{\partial q^i},
        \qquad
        \mathbf{e}^i=\nabla q^i.
    \]

    \[
        \mathbf{e}^i\cdot\mathbf{e}_j=\delta^i_j.
    \]

    \[
        \mathbf{A}=A^i\mathbf{e}_i=A_i\mathbf{e}^i.
    \]

    \[
        A^i=\mathbf{A}\cdot\mathbf{e}^i,
        \qquad
        A_i=\mathbf{A}\cdot\mathbf{e}_i.
    \]
\end{formulaboxnosplit}

\subsection{Metric Tensor}

The metric tensor records dot products between basis vectors. It is the object that turns coordinate changes into physical lengths and angles.

\begin{formulabox}[Metric Tensor and Raising/Lowering]
    \[
        g_{ij}=\mathbf{e}_i\cdot\mathbf{e}_j,
        \qquad
        g^{ij}=\mathbf{e}^i\cdot\mathbf{e}^j.
    \]

    \[
        g^{ik}g_{kj}=\delta^i_j.
    \]

    \[
        A_i=g_{ij}A^j,
        \qquad
        A^i=g^{ij}A_j.
    \]

    \[
        ds^2=g_{ij}\,dq^{i}dq^j.
    \]

    \[
        \mathbf{A}\cdot\mathbf{B}=g_{ij}A^{i}B^j=g^{ij}A_{i}B_j=A_{i}B^i=A^{i}B_i.
    \]
\end{formulabox}

\subsection{Orthogonal Coordinates as a Special Case}

If the coordinate system is orthogonal, the metric tensor is diagonal. This is why scale factors are enough to describe many engineering coordinate systems.

\begin{formulaboxnosplit}[Orthogonal Metric and Scale Factors]
    \[
        g_{ij}=h_i^2\delta_{ij}\quad\text{no sum on }i.
    \]

    \[
        g^{ij}=\frac{1}{h_i^2}\delta^{ij}\quad\text{no sum on }i.
    \]

    \[
        \mathbf{e}_i=h_i\hat{\mathbf{e}}_i,
        \qquad
        \mathbf{e}^i=\frac{1}{h_i}\hat{\mathbf{e}}_i.
    \]

    \[
        A_i=h_i A_{\text{phys},i},
        \qquad
        A^i=\frac{A_{\text{phys},i}}{h_i}.
    \]
\end{formulaboxnosplit}

\subsection{Jacobian, Metric and Volume}

For a coordinate transformation \(x^a=x^a(q^1,q^2,q^3)\), the Jacobian matrix connects small coordinate changes to small physical displacements.

\begin{formulabox}[Jacobian and Volume Element]
    \[
        J^a{}_i=\frac{\partial x^a}{\partial q^i}.
    \]

    \[
        g_{ij}=\frac{\partial x^a}{\partial q^i}\frac{\partial x^a}{\partial q^j}=J^a{}_{i}J^a{}_j.
    \]

    \[
        g=\det(g_{ij}),
        \qquad
        \sqrt{g}=\left|\det(J)\right|.
    \]

    \[
        dV=\sqrt{g}\,dq^1dq^2dq^3.
    \]

    For orthogonal coordinates,
    \[
        \sqrt{g}=h_1h_2h_3.
    \]
\end{formulabox}

\subsection{Gradient, Divergence and Laplacian in General Curvilinear Coordinates}

The following formulas use the metric determinant \(g=\det(g_{ij})\). They are especially useful when coordinates are not necessarily orthogonal.

\begin{formulabox}[General Curvilinear Formulae]
    \[
        \nabla f=(\partial_i f)\mathbf{e}^i.
    \]

    \[
        \nabla\cdot\mathbf{A}=\frac{1}{\sqrt{g}}\frac{\partial}{\partial q^i}\left(\sqrt{g}\,A^i\right).
    \]

    \[
        \nabla^2f=\frac{1}{\sqrt{g}}\frac{\partial}{\partial q^i}\left(\sqrt{g}\,g^{ij}\frac{\partial f}{\partial q^j}\right).
    \]

    For orthogonal coordinates,
    \[
        \nabla^2 f=\frac{1}{h_1h_2h_3}\sum_{i=1}^{3}\frac{\partial}{\partial q^i}\left(\frac{h_1h_2h_3}{h_i^2}\frac{\partial f}{\partial q^i}\right).
    \]
\end{formulabox}

\subsection{Cylindrical and Spherical Laplacians}

\begin{formulabox}[Common Laplacians]
    Cartesian:
    \[
        \nabla^2f=f_{xx}+f_{yy}+f_{zz}.
    \]

    Cylindrical:
    \[
        \nabla^2f=\frac{1}{\rho}\frac{\partial}{\partial\rho}\left(\rho\frac{\partial f}{\partial\rho}\right)+\frac{1}{\rho^2}\frac{\partial^2f}{\partial\phi^2}+\frac{\partial^2f}{\partial z^2}.
    \]

    Spherical:
    \[
        \nabla^2f=\frac{1}{r^2}\frac{\partial}{\partial r}\left(r^2\frac{\partial f}{\partial r}\right)
        +\frac{1}{r^2\sin\theta}\frac{\partial}{\partial\theta}\left(\sin\theta\frac{\partial f}{\partial\theta}\right)
        +\frac{1}{r^2\sin^2\theta}\frac{\partial^2f}{\partial\phi^2}.
    \]
\end{formulabox}

\begin{practicebox}[Quick Drills and Practice: Covariant/Contravariant Components and Metrics]
    \textbf{Level A:\,Conceptual recall}

        \begin{enumerate}
            \item Complete: \(\mathbf{e}_i=\blank[4cm]\).
            \item Complete: \(\mathbf{e}^i=\blank[3cm]\).
            \item Complete: \(\mathbf{e}^i\cdot\mathbf{e}_j=\blank[2cm]\).
            \item Complete: \(g_{ij}=\blank[4cm]\).
            \item Complete: \(g^{ik}g_{kj}=\blank[2cm]\).
            \item Complete: \(A_i=\blank[3cm]\).
            \item Complete: \(A^i=\blank[3cm]\).
            \item Complete: \(ds^2=\blank[4cm]\).
        \end{enumerate}


    \textbf{Level B:\,Metric computations}
    \begin{enumerate}
        \item Compute \(g_{ij}\) for cylindrical coordinates directly from \(\mathbf{r}(\rho,\phi,z)\).
        \item Compute \(g^{ij}\) for cylindrical coordinates.
        \item Compute \(g_{ij}\) for spherical coordinates directly from \(\mathbf{r}(r,\theta,\phi)\).
        \item Compute \(g^{ij}\) for spherical coordinates.
        \item Compute \(\sqrt{g}\) for Cartesian, cylindrical and spherical coordinates.
        \item Show that \(dV=\sqrt{g}\,dq^1dq^2dq^3\) gives the correct cylindrical and spherical volume elements.
    \end{enumerate}

    \textbf{Level C:\,Components}
    \begin{enumerate}
        \item In cylindrical coordinates, if the physical components are \((A_\rho,A_\phi,A_z)=(2,3,4)\), write the contravariant components \(A^i\).
        \item In spherical coordinates, if the physical components are \((A_r,A_\theta,A_\phi)=(1,2,3)\), write \(A^i\) and \(A_i\).
        \item Given \(g_{ij}=\diag(1,r^2,r^2\sin^2\theta)\), lower the index of \(A^i=(r,\sin\theta,1)\).
        \item Given \(g^{ij}=\diag(1,1/r^2,1/(r^2\sin^2\theta))\), raise the index of \(A_i=(r^2,r^2\theta,r^2\sin\theta)\).
    \end{enumerate}

    \textbf{Level D:\,Differential operators}
    \begin{enumerate}
        \item Derive the cylindrical divergence formula from \(\nabla\cdot\mathbf{A}=\frac1{\sqrt g}\partial_i(\sqrt g A^i)\).
        \item Derive the spherical divergence formula from \(\nabla\cdot\mathbf{A}=\frac1{\sqrt g}\partial_i(\sqrt g A^i)\).
        \item Derive the cylindrical scalar Laplacian from the general curvilinear formula.
        \item Derive the spherical scalar Laplacian from the general curvilinear formula.
        \item Compute \(\nabla^2f\) in cylindrical coordinates for \(f=\rho^2+z^2\).
        \item Compute \(\nabla^2f\) in spherical coordinates for \(f=r^2\).
    \end{enumerate}
\end{practicebox}

\ifdefined\OTCOMBINED%
\else
\end{document}
\fi

\ifdefined\OTCOMBINED%
\else
\documentclass{otscience}
% ============================================================
% Shared setup for OT Science modular workbook parts
% Place these files in the same folder as your fixed otscience.cls
% ============================================================

\makeatletter
\makeatother
\usepackage{multicol}
\usepackage{longtable}
\usepackage{makecell}
\usepackage{pdflscape}
\usepackage{bm}
\providecommand{\blank}[1][3cm]{\rule{#1}{0.4pt}}
\providecommand{\bigblank}{\rule{7cm}{0.4pt}}
\providecommand{\componentblank}{\blank[2.5cm]}
\providecommand{\R}{\mathbb{R}}
\providecommand{\C}{\mathbb{C}}
\providecommand{\ii}{\mathrm{i}}
\providecommand{\ee}{\mathrm{e}}
\providecommand{\dd}{\mathrm{d}}
\providecommand{\grad}{\nabla}
\providecommand{\divergence}{\nabla\!\cdot}
\providecommand{\curl}{\nabla\!\times}
\providecommand{\lap}{\nabla^2}
\providecommand{\ihat}{\hat{\mathbf{i}}}
\providecommand{\jhat}{\hat{\mathbf{j}}}
\providecommand{\khat}{\hat{\mathbf{k}}}
\providecommand{\er}{\hat{\mathbf{e}}_r}
\providecommand{\etheta}{\hat{\mathbf{e}}_\theta}
\providecommand{\ephi}{\hat{\mathbf{e}}_\phi}
\providecommand{\erho}{\hat{\boldsymbol{\rho}}}
\providecommand{\vA}{\mathbf{A}}
\providecommand{\vB}{\mathbf{B}}
\providecommand{\va}{\mathbf{a}}
\providecommand{\vb}{\mathbf{b}}
\providecommand{\vc}{\mathbf{c}}
\providecommand{\vx}{\mathbf{x}}
\providecommand{\vr}{\mathbf{r}}
\providecommand{\matA}{\mathbf{A}}
\providecommand{\matB}{\mathbf{B}}
\providecommand{\tr}{\operatorname{tr}}
\providecommand{\cof}{\operatorname{cof}}
\providecommand{\dev}{\operatorname{dev}}
\providecommand{\diag}{\operatorname{diag}}
\providecommand{\questionline}[1][5cm]{\rule{#1}{0.4pt}}
\setlength{\parindent}{0pt}
\setlength{\parskip}{4pt}
\renewcommand{\arraystretch}{1.25}

% Fallbacks in case the current otscience.cls has not defined nosplit boxes.
\makeatletter
\@ifundefined{formulaboxnosplit}{\newenvironment{formulaboxnosplit}[1][Formula]{\begin{formulabox}[#1]}{\end{formulabox}}}{}
\@ifundefined{definitionboxnosplit}{\newenvironment{definitionboxnosplit}[1][Definition]{\begin{definitionbox}[#1]}{\end{definitionbox}}}{}
\@ifundefined{summaryboxnosplit}{\newenvironment{summaryboxnosplit}[1][Summary]{\begin{summarybox}[#1]}{\end{summarybox}}}{}
\@ifundefined{warningboxnosplit}{\newenvironment{warningboxnosplit}[1][Warning]{\begin{warningbox}[#1]}{\end{warningbox}}}{}
\@ifundefined{exampleboxnosplit}{\newenvironment{exampleboxnosplit}[1][Example]{\begin{examplebox}[#1]}{\end{examplebox}}}{}
\makeatother

\newenvironment{practicebox}[1][Quick Drills and Practice]{\begin{examplebox}[#1]}{\end{examplebox}}
\newenvironment{practiceboxnosplit}[1][Quick Drills and Practice]{\begin{exampleboxnosplit}[#1]}{\end{exampleboxnosplit}}

\newcommand{\standalonetitle}[3]{%
    \title{\textbf{#1}\\\large #2}
    \author{OT Science Notes}
    \date{#3}
}

\standalonetitle{Symmetric and Skew-Symmetric Tensors, Invariants and Decompositions}{Part 5 of the OT Science Formula Completion Workbook}{\DTMtoday}
\begin{document}
\maketitle
\tableofcontents
\newpage
\fi

\section{Symm and Skew-Symm Tensors, Invariants  n' Decomps}

Second-order tensors appear naturally as stress tensors, strain tensors, inertia tensors, conductivity tensors, deformation gradients and vector gradients. Their symmetric and skew-symmetric parts often carry different physical meanings.

\subsection{Symmetric and Skew-Symmetric Tensors}

A tensor is symmetric if transposing it does not change it. A tensor is skew-symmetric if transposing it changes its sign. In three dimensions, a skew-symmetric tensor has only three independent components and can be represented by an axial vector.

\begin{formulaboxnosplit}[Symmetry Conditions]
    \[
        S_{ij}=S_{ji}\qquad\text{symmetry}.
    \]

    \[
        W_{ij}=-W_{ji}\qquad\text{skew-symmetry}.
    \]

    \[
        W_{ii}=0.
    \]

    \[
        A_{ij}=S_{ij}+W_{ij}.
    \]

    \[
        S_{ij}=\frac12(A_{ij}+A_{ji}),
        \qquad
        W_{ij}=\frac12(A_{ij}-A_{ji}).
    \]
\end{formulaboxnosplit}

\subsection{Connection with Vector Gradient}

If \(\mathbf{v}\) is a velocity field, then \(\nabla\mathbf{v}\) can be decomposed into symmetric and skew-symmetric parts. The symmetric part is related to strain rate, while the skew-symmetric part is related to local rotation.

\begin{formulabox}[Gradient Decomposition]
    \[
        L_{ij}=\partial_{j}v_i.
    \]

    \[
        D_{ij}=\frac12(L_{ij}+L_{ji})
        \qquad\text{symmetric rate-of-deformation tensor}.
    \]

    \[
        \Omega_{ij}=\frac12(L_{ij}-L_{ji})
        \qquad\text{skew spin tensor}.
    \]

    \[
        L_{ij}=D_{ij}+\Omega_{ij}.
    \]

    \[
        \omega_i={(\nabla\times\mathbf{v})}_i=\varepsilon_{ijk}\partial_{j}v_k.
    \]
\end{formulabox}

\subsection{Principal Invariants of a Second-Order Tensor}

Invariants do not change under a rotation of coordinates. They describe intrinsic properties of the tensor.

\begin{formulaboxnosplit}[Principal Invariants]
    For a \(3\times3\) tensor \(\mathbf{A}\),
    \[
        I_1=\tr(A)=A_{ii}.
    \]

    \[
        I_2=\frac12\left[{(\tr A)}^2-\tr(A^2)\right].
    \]

    \[
        I_3=\det(A).
    \]

    The characteristic polynomial may be written as
    \[
        \det(\lambda I-A)=\lambda^3-I_1\lambda^2+I_2\lambda-I_3.
    \]
    Equivalently,
    \[
        \det(A-\lambda I)=-\lambda^3+I_1\lambda^2-I_2\lambda+I_3.
    \]
\end{formulaboxnosplit}

\subsection{Deviatoric and Spherical Parts}

Many physical tensors can be decomposed into a mean or spherical part and a deviatoric part. In stress analysis, the spherical part is associated with hydrostatic pressure, while the deviatoric part is associated with shape distortion.

\begin{formulabox}[Deviatoric Decomposition]
    \[
        \operatorname{sph}{(A)}_{ij}=\frac13A_{kk}\delta_{ij}.
    \]

    \[
        A'_{ij}=\dev{(A)}_{ij}=A_{ij}-\frac13A_{kk}\delta_{ij}.
    \]

    \[
        A_{ij}=\frac13A_{kk}\delta_{ij}+A'_{ij}.
    \]

    \[
        A'_{kk}=0.
    \]

    For a stress tensor \(\sigma_{ij}\),
    \[
        p=-\frac13\sigma_{kk},
        \qquad
        s_{ij}=\sigma_{ij}-\frac13\sigma_{kk}\delta_{ij}.
    \]
\end{formulabox}

\subsection{Invariants of Deviatoric Tensors}

For a deviatoric tensor \(A'\), the first invariant is zero. The remaining invariants are important in mechanics.

\begin{formulabox}[Deviatoric Invariants]
    \[
        J_1=\tr(A')=0.
    \]

    \[
        J_2=\frac12A'_{ij}A'_{ij}.
    \]

    \[
        J_3=\det(A').
    \]

    For stress deviator \(s_{ij}\), one common convention is
    \[
        J_2=\frac12s_{ij}s_{ij}.
    \]
\end{formulabox}

\begin{exampleboxnosplit}[Worked Example: Tensor Decomposition]
    Let
    \[
        A=\begin{pmatrix}
            2 & 1 & 3  \\
            4 & 0 & -1 \\
            5 & 2 & 1
        \end{pmatrix}.
    \]
    Then
    \[
        S=\frac12(A+A^T),
        \qquad
        W=\frac12(A-A^T).
    \]
    Compute
    \[
        S=\begin{pmatrix}
            2   & 5/2 & 4   \\
            5/2 & 0   & 1/2 \\
            4   & 1/2 & 1
        \end{pmatrix},
        \qquad
        W=\begin{pmatrix}
            0   & -3/2 & -1   \\
            3/2 & 0    & -3/2 \\
            1   & 3/2  & 0
        \end{pmatrix}.
    \]
    Check that \(S^T=S\), \(W^T=-W\), and \(A=S+W\).
\end{exampleboxnosplit}

\begin{practicebox}[Quick Drills and Practice: Tensor Decompositions and Invariants]
    \textbf{Level A:\,Symmetry checks}
    \begin{multicols}{2}
        \begin{enumerate}
            \item Is \(\begin{pmatrix}1&2\\2&3\end{pmatrix}\) symmetric?
            \item Is \(\begin{pmatrix}0&-4\\4&0\end{pmatrix}\) skew-symmetric?
            \item Complete: \(S_{ij}-S_{ji}=\blank[2cm]\).
            \item Complete: \(W_{ij}+W_{ji}=\blank[2cm]\).
            \item Complete: \(W_{ii}=\blank[2cm]\).
            \item Complete: \(A_{ij}=\blank[3cm]\).
            \item Complete: \(S_{ij}=\blank[3cm]\).
            \item Complete: \(W_{ij}=\blank[3cm]\).
        \end{enumerate}
    \end{multicols}

    \textbf{Level B:\,Decomposition}
    \begin{enumerate}
        \item Decompose \(A=\begin{pmatrix}1&2&3\\0&4&5\\1&0&6\end{pmatrix}\) into symmetric and skew-symmetric parts.
        \item Decompose \(A=\begin{pmatrix}0&-1&2\\1&0&-3\\-2&3&0\end{pmatrix}\). What do you notice?
        \item For \(\mathbf{v}=(-y,x,0)\), compute \(\nabla\mathbf{v}\), its symmetric part and its skew part.
        \item For \(\mathbf{v}=(ax,by,cz)\), compute \(\nabla\mathbf{v}\), its trace and its symmetric/skew parts.
    \end{enumerate}

    \textbf{Level C:\,Invariants}
    \begin{enumerate}
        \item Compute \(I_1,I_2,I_3\) for \(A=\diag(1,2,3)\).
        \item Compute \(I_1,I_2,I_3\) for \(A=\diag(a,b,c)\).
        \item Compute \(I_1,I_2,I_3\) for \(A=\begin{pmatrix}2&1&0\\1&2&0\\0&0&3\end{pmatrix}\).
        \item Write the characteristic polynomial of a \(3\times3\) tensor using \(I_1,I_2,I_3\).
        \item If a tensor has eigenvalues \(\lambda_1,\lambda_2,\lambda_3\), express \(I_1,I_2,I_3\) in terms of the eigenvalues.
    \end{enumerate}

    \textbf{Level D:\,Deviatoric tensors}
    \begin{enumerate}
        \item Find the deviatoric part of \(A=\diag(4,1,1)\).
        \item Find the deviatoric part of \(A=\diag(6,3,0)\).
        \item Prove that \(A'_{kk}=0\).
        \item For \(A=\diag(4,1,1)\), compute \(J_2=\frac12A'_{ij}A'_{ij}\).
        \item For \(\sigma=\diag(100,50,25)\), compute the hydrostatic part and deviatoric part.
    \end{enumerate}

    \textbf{Level E:\,Challenge}
    \begin{enumerate}
        \item Show that the trace of a skew-symmetric tensor is zero.
        \item Show that \(x_{i}W_{ij}x_j=0\) for any skew-symmetric tensor \(W_{ij}\).
        \item Show that \(x_{i}S_{ij}x_j\) depends only on the symmetric part of a tensor.
        \item If \(A\) is symmetric, explain why its eigenvalues are real in ordinary Euclidean space.
        \item Explain why invariants are useful when comparing the same tensor in two different coordinate frames.
    \end{enumerate}
\end{practicebox}

\ifdefined\OTCOMBINED%
\else
\end{document}
\fi

\ifdefined\OTCOMBINED%
\else
\documentclass{otscience}
% ============================================================
% Shared setup for OT Science modular workbook parts
% Place these files in the same folder as your fixed otscience.cls
% ============================================================

\makeatletter
\makeatother
\usepackage{multicol}
\usepackage{longtable}
\usepackage{makecell}
\usepackage{pdflscape}
\usepackage{bm}
\providecommand{\blank}[1][3cm]{\rule{#1}{0.4pt}}
\providecommand{\bigblank}{\rule{7cm}{0.4pt}}
\providecommand{\componentblank}{\blank[2.5cm]}
\providecommand{\R}{\mathbb{R}}
\providecommand{\C}{\mathbb{C}}
\providecommand{\ii}{\mathrm{i}}
\providecommand{\ee}{\mathrm{e}}
\providecommand{\dd}{\mathrm{d}}
\providecommand{\grad}{\nabla}
\providecommand{\divergence}{\nabla\!\cdot}
\providecommand{\curl}{\nabla\!\times}
\providecommand{\lap}{\nabla^2}
\providecommand{\ihat}{\hat{\mathbf{i}}}
\providecommand{\jhat}{\hat{\mathbf{j}}}
\providecommand{\khat}{\hat{\mathbf{k}}}
\providecommand{\er}{\hat{\mathbf{e}}_r}
\providecommand{\etheta}{\hat{\mathbf{e}}_\theta}
\providecommand{\ephi}{\hat{\mathbf{e}}_\phi}
\providecommand{\erho}{\hat{\boldsymbol{\rho}}}
\providecommand{\vA}{\mathbf{A}}
\providecommand{\vB}{\mathbf{B}}
\providecommand{\va}{\mathbf{a}}
\providecommand{\vb}{\mathbf{b}}
\providecommand{\vc}{\mathbf{c}}
\providecommand{\vx}{\mathbf{x}}
\providecommand{\vr}{\mathbf{r}}
\providecommand{\matA}{\mathbf{A}}
\providecommand{\matB}{\mathbf{B}}
\providecommand{\tr}{\operatorname{tr}}
\providecommand{\cof}{\operatorname{cof}}
\providecommand{\dev}{\operatorname{dev}}
\providecommand{\diag}{\operatorname{diag}}
\providecommand{\questionline}[1][5cm]{\rule{#1}{0.4pt}}
\setlength{\parindent}{0pt}
\setlength{\parskip}{4pt}
\renewcommand{\arraystretch}{1.25}

% Fallbacks in case the current otscience.cls has not defined nosplit boxes.
\makeatletter
\@ifundefined{formulaboxnosplit}{\newenvironment{formulaboxnosplit}[1][Formula]{\begin{formulabox}[#1]}{\end{formulabox}}}{}
\@ifundefined{definitionboxnosplit}{\newenvironment{definitionboxnosplit}[1][Definition]{\begin{definitionbox}[#1]}{\end{definitionbox}}}{}
\@ifundefined{summaryboxnosplit}{\newenvironment{summaryboxnosplit}[1][Summary]{\begin{summarybox}[#1]}{\end{summarybox}}}{}
\@ifundefined{warningboxnosplit}{\newenvironment{warningboxnosplit}[1][Warning]{\begin{warningbox}[#1]}{\end{warningbox}}}{}
\@ifundefined{exampleboxnosplit}{\newenvironment{exampleboxnosplit}[1][Example]{\begin{examplebox}[#1]}{\end{examplebox}}}{}
\makeatother

\newenvironment{practicebox}[1][Quick Drills and Practice]{\begin{examplebox}[#1]}{\end{examplebox}}
\newenvironment{practiceboxnosplit}[1][Quick Drills and Practice]{\begin{exampleboxnosplit}[#1]}{\end{exampleboxnosplit}}

\newcommand{\standalonetitle}[3]{%
    \title{\textbf{#1}\\\large #2}
    \author{OT Science Notes}
    \date{#3}
}

\standalonetitle{Complex Numbers, Circular Functions and Hyperbolic Functions}{Part 6 of the OT Science Formula Completion Workbook}{\DTMtoday}
\begin{document}
\maketitle
\tableofcontents
\newpage
\fi

\section{Complex Numbers, Circular Functions and Hyperbolic Functions}

Complex numbers connect algebra, geometry, oscillations, rotations, exponentials and hyperbolic functions. The core bridge is Euler's formula. Once that is understood, trigonometric and hyperbolic identities become different faces of the same exponential structure.

\subsection{Basic Complex Algebra}

A complex number has a real part and an imaginary part:
\[
z=x+\ii y,
\qquad \ii^2=-1.
\]
The complex plane treats \(x\) as the horizontal coordinate and \(y\) as the vertical coordinate.

\begin{formulaboxnosplit}[Complex Number Basics]
\[
z=x+\ii y,
\qquad
\Re(z)=x,
\qquad
\Im(z)=y.
\]

\[
\overline{z}=x-\ii y.
\]

\[
|z|=\sqrt{x^2+y^2},
\qquad
\arg(z)=\theta.
\]

\[
z\overline{z}=|z|^2.
\]

\[
\frac{1}{z}=\frac{\overline{z}}{|z|^2}=\frac{x-\ii y}{x^2+y^2},\qquad z\neq0.
\]
\end{formulaboxnosplit}

\subsection{Polar and Exponential Forms}

The polar form is often the best way to multiply, divide, raise to powers and extract roots.

\begin{formulabox}[Euler, Polar Form and De Moivre]
\[
\ee^{\ii\theta}=\cos\theta+\ii\sin\theta.
\]

\[
z=r(\cos\theta+\ii\sin\theta)=r\ee^{\ii\theta}.
\]

\[
r=|z|,
\qquad
\theta=\arg(z).
\]

\[
z_1z_2=r_1r_2\ee^{\ii(\theta_1+\theta_2)}.
\]

\[
\frac{z_1}{z_2}=\frac{r_1}{r_2}\ee^{\ii(\theta_1-\theta_2)},\qquad z_2\neq0.
\]

\[
{(r\ee^{\ii\theta})}^n=r^n\ee^{\ii n\theta}.
\]

\[
{(\cos\theta+\ii\sin\theta)}^n=\cos(n\theta)+\ii\sin(n\theta).
\]
\end{formulabox}

\subsection{Roots of Complex Numbers}

If \(z^n=R\ee^{\ii\Theta}\), then the \(n\) roots are evenly spaced around a circle in the complex plane.

\begin{formulaboxnosplit}[Complex Roots]
\[
z_k=R^{1/n}\ee^{\ii(\Theta+2\pi k)/n},
\qquad k=0,1,2,\dots,n-1.
\]

Equivalently,
\[
z_k=R^{1/n}\left[\cos\left(\frac{\Theta+2\pi k}{n}\right)+\ii\sin\left(\frac{\Theta+2\pi k}{n}\right)\right].
\]
\end{formulaboxnosplit}

\subsection{Complex Exponentials and Circular Functions}

The ordinary circular functions can be written using complex exponentials:

\begin{formulabox}[Circular Functions from Exponentials]
\[
\cos x=\frac{\ee^{\ii x}+\ee^{-\ii x}}{2}.
\]

\[
\sin x=\frac{\ee^{\ii x}-\ee^{-\ii x}}{2\ii}.
\]

\[
\tan x=\frac{\sin x}{\cos x}.
\]

\[
\cos^2x+\sin^2x=1.
\]

\[
\ee^{\ii x}\ee^{-\ii x}=1.
\]
\end{formulabox}

\subsection{Hyperbolic Functions}

Hyperbolic functions are built from real exponentials. They resemble circular functions, but their signs differ in crucial places.

\begin{formulabox}[Hyperbolic Definitions and Identities]
\[
\sinh x=\frac{\ee^x-\ee^{-x}}{2}.
\]

\[
\cosh x=\frac{\ee^x+\ee^{-x}}{2}.
\]

\[
\tanh x=\frac{\sinh x}{\cosh x}.
\]

\[
\cosh^2x-\sinh^2x=1.
\]

\[
\sinh(2x)=2\sinh x\cosh x.
\]

\[
\cosh(2x)=\cosh^2x+\sinh^2x=2\cosh^2x-1=1+2\sinh^2x.
\]

\[
\frac{d}{dx}\sinh x=\cosh x,
\qquad
\frac{d}{dx}\cosh x=\sinh x,
\qquad
\frac{d}{dx}\tanh x=\operatorname{sech}^2x.
\]
\end{formulabox}

\subsection{Circular---Hyperbolic Connections}

Replacing a real argument by an imaginary one converts circular functions into hyperbolic functions and vice versa.

\begin{formulaboxnosplit}[Complex Arguments and Hyperbolic Connections]
\[
\cos(\ii x)=\cosh x.
\]

\[
\sin(\ii x)=\ii\sinh x.
\]

\[
\tan(\ii x)=\ii\tanh x.
\]

\[
\cosh(\ii x)=\cos x.
\]

\[
\sinh(\ii x)=\ii\sin x.
\]

\[
\tanh(\ii x)=\ii\tan x.
\]
\end{formulaboxnosplit}

\subsection{Inverse Hyperbolic Functions}

Inverse hyperbolic functions can be written using logarithms. These are important in integration and differential equations.

\begin{formulabox}[Inverse Hyperbolic Forms]
\[
\sinh^{-1}x=\ln(x+\sqrt{x^2+1}).
\]

\[
\cosh^{-1}x=\ln(x+\sqrt{x^2-1}),\qquad x\ge1.
\]

\[
\tanh^{-1}x=\frac12\ln\left(\frac{1+x}{1-x}\right),\qquad |x|<1.
\]
\end{formulabox}

\begin{practicebox}[Quick Drills and Practice: Complex and Hyperbolic Functions]
\textbf{Level A:\,Formula recall}
\begin{multicols}{2}
\begin{enumerate}
\item Complete: \(\ee^{\ii\theta}=\blank[4cm]\).
\item Complete: \(z\overline{z}=\blank[2.5cm]\).
\item Complete: \(1/(x+\ii y)=\blank[4cm]\).
\item Complete: \(\cos x=\blank[4cm]\).
\item Complete: \(\sin x=\blank[4cm]\).
\item Complete: \(\sinh x=\blank[4cm]\).
\item Complete: \(\cosh x=\blank[4cm]\).
\item Complete: \(\cosh^2x-\sinh^2x=\blank[2cm]\).
\item Complete: \(\cos(\ii x)=\blank[3cm]\).
\item Complete: \(\sin(\ii x)=\blank[3cm]\).
\end{enumerate}
\end{multicols}

\textbf{Level B:\,Complex algebra}
\begin{enumerate}
\item Compute \((3+2\ii)+(1-5\ii)\).
\item Compute \((2+\ii)(3-4\ii)\).
\item Compute \({(1+\ii)}^4\) using ordinary algebra.
\item Write \(1+\ii\) in polar and exponential form.
\item Write \(-1+\sqrt3\ii\) in polar and exponential form.
\item Compute \({(1+\ii)}^8\) using polar form.
\item Compute \(\frac{2+3\ii}{1-\ii}\).
\end{enumerate}

\textbf{Level C:\,Powers and roots}
\begin{enumerate}
\item Find all cube roots of \(8\).
\item Find all fourth roots of \(1\).
\item Find all fourth roots of \(-16\).
\item Solve \(z^3=-8\ii\).
\item Sketch the roots of \(z^5=1\) on the complex plane.
\end{enumerate}

\textbf{Level D:\,Hyperbolic connections}
\begin{enumerate}
\item Prove \(\cos(\ii x)=\cosh x\) using exponential definitions.
\item Prove \(\sin(\ii x)=\ii\sinh x\).
\item Prove \(\cosh(\ii x)=\cos x\).
\item Prove \(\sinh(\ii x)=\ii\sin x\).
\item Prove \(\cosh^2x-\sinh^2x=1\).
\item Derive \(\sinh(2x)=2\sinh x\cosh x\).
\item Derive \(\cosh(2x)=\cosh^2x+\sinh^2x\).
\end{enumerate}

\textbf{Level E:\,Challenge}
\begin{enumerate}
\item Express \(\cos(x+\ii y)\) in terms of \(\cos x,\sin x,\cosh y,\sinh y\).
\item Express \(\sin(x+\ii y)\) in terms of \(\cos x,\sin x,\cosh y,\sinh y\).
\item Solve \(\ee^z=-1\) for complex \(z\).
\item Solve \(\cos z=0\) for complex \(z\).
\item Show that multiplication by \(\ee^{\ii\theta}\) rotates a complex number by angle \(\theta\).
\end{enumerate}
\end{practicebox}

\ifdefined\OTCOMBINED%
\else
\end{document}
\fi

\ifdefined\OTCOMBINED%
\else
\documentclass{otscience}
% ============================================================
% Shared setup for OT Science modular workbook parts
% Place these files in the same folder as your fixed otscience.cls
% ============================================================

\makeatletter
\makeatother
\usepackage{multicol}
\usepackage{longtable}
\usepackage{makecell}
\usepackage{pdflscape}
\usepackage{bm}
\providecommand{\blank}[1][3cm]{\rule{#1}{0.4pt}}
\providecommand{\bigblank}{\rule{7cm}{0.4pt}}
\providecommand{\componentblank}{\blank[2.5cm]}
\providecommand{\R}{\mathbb{R}}
\providecommand{\C}{\mathbb{C}}
\providecommand{\ii}{\mathrm{i}}
\providecommand{\ee}{\mathrm{e}}
\providecommand{\dd}{\mathrm{d}}
\providecommand{\grad}{\nabla}
\providecommand{\divergence}{\nabla\!\cdot}
\providecommand{\curl}{\nabla\!\times}
\providecommand{\lap}{\nabla^2}
\providecommand{\ihat}{\hat{\mathbf{i}}}
\providecommand{\jhat}{\hat{\mathbf{j}}}
\providecommand{\khat}{\hat{\mathbf{k}}}
\providecommand{\er}{\hat{\mathbf{e}}_r}
\providecommand{\etheta}{\hat{\mathbf{e}}_\theta}
\providecommand{\ephi}{\hat{\mathbf{e}}_\phi}
\providecommand{\erho}{\hat{\boldsymbol{\rho}}}
\providecommand{\vA}{\mathbf{A}}
\providecommand{\vB}{\mathbf{B}}
\providecommand{\va}{\mathbf{a}}
\providecommand{\vb}{\mathbf{b}}
\providecommand{\vc}{\mathbf{c}}
\providecommand{\vx}{\mathbf{x}}
\providecommand{\vr}{\mathbf{r}}
\providecommand{\matA}{\mathbf{A}}
\providecommand{\matB}{\mathbf{B}}
\providecommand{\tr}{\operatorname{tr}}
\providecommand{\cof}{\operatorname{cof}}
\providecommand{\dev}{\operatorname{dev}}
\providecommand{\diag}{\operatorname{diag}}
\providecommand{\questionline}[1][5cm]{\rule{#1}{0.4pt}}
\setlength{\parindent}{0pt}
\setlength{\parskip}{4pt}
\renewcommand{\arraystretch}{1.25}

% Fallbacks in case the current otscience.cls has not defined nosplit boxes.
\makeatletter
\@ifundefined{formulaboxnosplit}{\newenvironment{formulaboxnosplit}[1][Formula]{\begin{formulabox}[#1]}{\end{formulabox}}}{}
\@ifundefined{definitionboxnosplit}{\newenvironment{definitionboxnosplit}[1][Definition]{\begin{definitionbox}[#1]}{\end{definitionbox}}}{}
\@ifundefined{summaryboxnosplit}{\newenvironment{summaryboxnosplit}[1][Summary]{\begin{summarybox}[#1]}{\end{summarybox}}}{}
\@ifundefined{warningboxnosplit}{\newenvironment{warningboxnosplit}[1][Warning]{\begin{warningbox}[#1]}{\end{warningbox}}}{}
\@ifundefined{exampleboxnosplit}{\newenvironment{exampleboxnosplit}[1][Example]{\begin{examplebox}[#1]}{\end{examplebox}}}{}
\makeatother

\newenvironment{practicebox}[1][Quick Drills and Practice]{\begin{examplebox}[#1]}{\end{examplebox}}
\newenvironment{practiceboxnosplit}[1][Quick Drills and Practice]{\begin{exampleboxnosplit}[#1]}{\end{exampleboxnosplit}}

\newcommand{\standalonetitle}[3]{%
    \title{\textbf{#1}\\\large #2}
    \author{OT Science Notes}
    \date{#3}
}

\standalonetitle{Final Mixed Practice Bank and Formula Recall}{Part 7 of the OT Science Formula Completion Workbook}{\DTMtoday}
\begin{document}
\maketitle
\tableofcontents
\newpage
\fi

% 07_final_mixed_practice_bank.tex
\section{Final Mixed Practice Bank and Formula Recall}
\sectionmark{Mixed Practice Bank}

This final part is designed for spaced revision. It deliberately mixes recall, computation, proof, interpretation and engineering-style applications. Use it after completing the earlier parts, or use it diagnostically to find weak areas.

\subsection{Formula Recall Page}

\begin{summarybox}[Final Formula Recall]
Fill these from memory before checking earlier sections.
\begin{multicols}{2}
\begin{enumerate}
\item \(\nabla f=\blank[5cm]\)
\item \(\nabla\cdot\mathbf{A}=\blank[5cm]\)
\item \(\nabla\times\mathbf{A}=\blank[5cm]\)
\item \(D_{\hat{\mathbf{u}}}f=\blank[4cm]\)
\item \(\nabla^2 f=\blank[5cm]\)
\item \({(\nabla\mathbf{A})}_{ij}=\blank[4cm]\)
\item \(\delta_{ij}a_j=\blank[3cm]\)
\item \(\varepsilon_{ijk}\varepsilon_{ilm}=\blank[4cm]\)
\item \(S_{ij}=\blank[4cm]\)
\item \(W_{ij}=\blank[4cm]\)
\item \(I_1=\blank[3cm]\)
\item \(I_2=\blank[5cm]\)
\item \(I_3=\blank[3cm]\)
\item \(A'_{ij}=\blank[5cm]\)
\item \(\mathbf{g}_i=\blank[4cm]\)
\item \(\mathbf{g}^i\cdot\mathbf{g}_j=\blank[3cm]\)
\item \(A_i=\blank[4cm]\)
\item \(A^i=\blank[4cm]\)
\item \(ds^2=\blank[5cm]\)
\item \(dV=\blank[5cm]\)
\item \(dV_{\text{cyl}}=\blank[4cm]\)
\item \(dV_{\text{sph}}=\blank[4cm]\)
\item \(\ee^{\ii\theta}=\blank[4cm]\)
\item \(\cos(\ii x)=\blank[3cm]\)
\item \(\sin(\ii x)=\blank[3cm]\)
\item \(\cosh^2x-\sinh^2x=\blank[2cm]\)
\end{enumerate}
\end{multicols}
\end{summarybox}

\subsection{Part I: Vector Calculus Computation}

\begin{practicebox}[Final Practice I:\,Grad, Div, Curl and Directional Derivatives]
\begin{enumerate}
\item Compute \(\nabla f\) for \(f=x^3y^2+z\cos x\).
\item Compute \(\nabla f\) for \(f=\ln(x^2+y^2+z^2)\), away from the origin.
\item Compute \(\nabla^2 f\) for \(f=x^2+y^2+z^2\).
\item Compute \(\nabla^2 f\) for \(f=1/r\), where \(r=\sqrt{x^2+y^2+z^2}\), away from the origin.
\item Compute \(\nabla\cdot\mathbf{A}\) for \(\mathbf{A}=x^2\ihat+y^2\jhat+z^2\khat\).
\item Compute \(\nabla\cdot\mathbf{A}\) for \(\mathbf{A}=yz\ihat+xz\jhat+xy\khat\).
\item Compute \(\nabla\times\mathbf{A}\) for \(\mathbf{A}=xy\ihat+yz\jhat+zx\khat\).
\item Compute \(\nabla\times\mathbf{A}\) for \(\mathbf{A}=(-y/(x^2+y^2))\ihat+(x/(x^2+y^2))\jhat\), away from the \(z\)-axis.
\item For \(f=x^2y+yz^2\), compute the directional derivative at \((1,1,2)\) in direction \((2,-1,2)\).
\item Find the direction of steepest increase of \(f=xy+yz+zx\) at \((1,2,3)\).
\item Compute \(\nabla\mathbf{A}\) for \(\mathbf{A}=x^2\ihat+xy\jhat+xz\khat\).
\item For the vector field in Question 11, find the symmetric and skew-symmetric parts of \(\nabla\mathbf{A}\).
\item Verify \(\nabla\cdot(\nabla\times\mathbf{A})=0\) for \(\mathbf{A}=z\ihat+x\jhat+y\khat\).
\item Verify \(\nabla\times(\nabla f)=0\) for \(f=xyz+x^2+y^2\).
\item Evaluate \(\nabla\times(\nabla\times\mathbf{A})\) directly and using the identity for \(\mathbf{A}=x^2\ihat+y^2\jhat+z^2\khat\).
\end{enumerate}
\end{practicebox}

\subsection{Part II: Coordinate Systems and Differential Operators}

\begin{practicebox}[Final Practice II:\,Coordinates, Scale Factors and Operators]
\begin{enumerate}
\item Convert \((x,y,z)=(-1,\sqrt3,2)\) to cylindrical coordinates using \(\operatorname{atan2}\).
\item Convert \((x,y,z)=(1,1,\sqrt2)\) to spherical coordinates.
\item Convert \((\rho,\phi,z)=(3,5\pi/6,-2)\) to Cartesian coordinates.
\item Convert \((r,\theta,\phi)=(2,\pi/3,\pi/4)\) to Cartesian coordinates.
\item Derive the cylindrical basis vectors \(\erho\) and \(\hat{\boldsymbol{\phi}}\) in terms of \(\ihat,\jhat\).
\item Derive the spherical basis vectors \(\er,\etheta,\ephi\) in terms of \(\ihat,\jhat,\khat\).
\item Derive \(h_\rho,h_\phi,h_z\) from the cylindrical position vector.
\item Derive \(h_r,h_\theta,h_\phi\) from the spherical position vector.
\item Write \([g_{ij}]\) and \([g^{ij}]\) for cylindrical coordinates.
\item Write \([g_{ij}]\) and \([g^{ij}]\) for spherical coordinates.
\item Compute \(\nabla f\) in cylindrical coordinates for \(f=\rho z\cos\phi\).
\item Compute \(\nabla f\) in spherical coordinates for \(f=r^3\sin\theta\cos\phi\).
\item Compute \(\nabla\cdot\mathbf{A}\) in cylindrical coordinates for \(\mathbf{A}=\rho\erho+\rho^2\hat{\boldsymbol{\phi}}+z\khat\).
\item Compute \(\nabla\cdot\mathbf{A}\) in spherical coordinates for \(\mathbf{A}=r\er+\sin\theta\,\etheta+\cos\phi\,\ephi\).
\item Compute \(\nabla\times\mathbf{A}\) in cylindrical coordinates for \(\mathbf{A}=0\erho+\rho\hat{\boldsymbol{\phi}}+z\khat\).
\item Compute \(\nabla\times\mathbf{A}\) in spherical coordinates for \(\mathbf{A}=0\er+0\etheta+r\sin\theta\ephi\).
\item Derive the cylindrical scalar Laplacian from the general orthogonal coordinate formula.
\item Derive the spherical scalar Laplacian from the general curvilinear formula.
\end{enumerate}
\end{practicebox}

\subsection{Part III: Index Notation and Tensor Identities}

\begin{practicebox}[Final Practice III:\,Index Notation, Delta and Levi-Civita]
\begin{enumerate}
\item Expand \(A_{ij}B_{jk}\) for the component \((i,k)=(1,2)\).
\item Simplify \(\delta_{ij}\delta_{jk}a_k\).
\item Simplify \(\delta_{ij}A_{ik}B_{jk}\).
\item Simplify \(\varepsilon_{ijk}\varepsilon_{ilm}a_{j}b_{k}c_{l}d_m\).
\item Show that \(\varepsilon_{ijk}a_{j}b_k\) gives the cross product.
\item Prove \(\mathbf{a}\times\mathbf{b}=-\mathbf{b}\times\mathbf{a}\).
\item Prove \((\mathbf{a}\times\mathbf{b})\cdot\mathbf{c}=\mathbf{a}\cdot(\mathbf{b}\times\mathbf{c})\).
\item Prove the vector triple product identity using \(\varepsilon_{ijk}\).
\item Prove \(\nabla\times(\nabla f)=0\) using indices.
\item Prove \(\nabla\cdot(\nabla\times\mathbf{A})=0\) using indices.
\item Prove the curl-curl identity using indices.
\item Compute \(\partial_i(x_{j}x_j)\).
\item Compute \(\partial_j(x_{i}x_j)\) in three dimensions.
\item Compute \(\partial_i(r^n)\), where \(r={(x_{j}x_j)}^{1/2}\).
\item Compute \(\partial_i(x_i/r^3)\) for \(r\neq0\).
\end{enumerate}
\end{practicebox}

\subsection{Part IV: Tensor Decompositions and Invariants}

\begin{practicebox}[Final Practice IV:\,Symmetry, Skew-Symmetry and Invariants]
\begin{enumerate}
\item Decompose \(A=\begin{pmatrix}2&1&0\\3&4&5\\1&-2&6\end{pmatrix}\) into symmetric and skew parts.
\item Check whether \(A=\begin{pmatrix}0&-2&1\\2&0&-3\\-1&3&0\end{pmatrix}\) is skew-symmetric.
\item Compute \(I_1,I_2,I_3\) for \(A=\diag(2,4,6)\).
\item Compute \(I_1,I_2,I_3\) for \(A=\begin{pmatrix}1&1&0\\1&1&0\\0&0&2\end{pmatrix}\).
\item Find the characteristic polynomial of \(A=\diag(a,b,c)\).
\item Find the deviatoric part of \(A=\diag(5,2,2)\).
\item Find the hydrostatic and deviatoric parts of \(\sigma=\diag(120,60,30)\).
\item Prove that the deviatoric tensor has zero trace.
\item Prove that \(x_{i}W_{ij}x_j=0\) for skew-symmetric \(W\).
\item Show that \(x_{i}A_{ij}x_j\) depends only on the symmetric part of \(A\).
\item For \(\mathbf{v}=(-\omega y,\omega x,0)\), compute the velocity gradient and split it into symmetric and skew parts.
\item For \(\mathbf{v}=(ax,by,cz)\), compute the invariants of \(\nabla\mathbf{v}\).
\end{enumerate}
\end{practicebox}

\subsection{Part V: Covariant and Contravariant Components}

\begin{practicebox}[Final Practice V:\,Metrics, Raising and Lowering]
\begin{enumerate}
\item Explain the difference between \(\mathbf{g}_i\) and \(\mathbf{g}^i\).
\item Prove \(\mathbf{g}^i\cdot\mathbf{g}_j=\delta^i_j\) for an orthogonal coordinate system.
\item In cylindrical coordinates, compute \(\mathbf{g}_\rho,\mathbf{g}_\phi,\mathbf{g}_z\).
\item In cylindrical coordinates, compute \(\mathbf{g}^\rho,\mathbf{g}^\phi,\mathbf{g}^z\).
\item In spherical coordinates, compute \(\mathbf{g}_r,\mathbf{g}_\theta,\mathbf{g}_\phi\).
\item In spherical coordinates, compute \(\mathbf{g}^r,\mathbf{g}^\theta,\mathbf{g}^\phi\).
\item Given cylindrical physical components \((A_\rho,A_\phi,A_z)=(3,4,5)\), find \(A^i\) and \(A_i\).
\item Given spherical physical components \((A_r,A_\theta,A_\phi)=(2,3,4)\), find \(A^i\) and \(A_i\).
\item Prove \(\mathbf{A}\cdot\mathbf{B}=g_{ij}A^{i}B^j\).
\item Derive \(\nabla^2f=\frac1{\sqrt g}\partial_i(\sqrt g g^{ij}\partial_{j}f)\) for the cylindrical case.
\end{enumerate}
\end{practicebox}

\subsection{Part VI: Complex Numbers and Hyperbolic Functions}

\begin{practicebox}[Final Practice VI:\,Complex and Hyperbolic Problems]
\begin{enumerate}
\item Write \(z=-1+\ii\) in polar and exponential form.
\item Write \(z=-\sqrt3-\ii\) in polar and exponential form.
\item Compute \({(\sqrt3+\ii)}^6\) using polar form.
\item Find all roots of \(z^4=16\).
\item Find all roots of \(z^3=-8\).
\item Solve \(z^5=32\ee^{\ii\pi/2}\).
\item Prove Euler's formula from the power series of \(\ee^x\), \(\sin x\), and \(\cos x\).
\item Derive De Moivre's theorem.
\item Express \(\cos(x+\ii y)\) in real and imaginary parts.
\item Express \(\sin(x+\ii y)\) in real and imaginary parts.
\item Prove \(\cos(\ii x)=\cosh x\).
\item Prove \(\sin(\ii x)=\ii\sinh x\).
\item Prove \(\tanh(\ii x)=\ii\tan x\).
\item Derive \(\sinh^{-1}x=\ln(x+\sqrt{x^2+1})\).
\item Derive \(\tanh^{-1}x=\frac12\ln\left(\frac{1+x}{1-x}\right)\).
\end{enumerate}
\end{practicebox}

\subsection{Part VII: Mixed Challenge Problems}

\begin{practicebox}[Final Practice VII: Mixed Challenge Set]
\begin{enumerate}
\item A scalar field is given by \(T=1/r\). Compute \(\nabla T\), \(\nabla^2T\) for \(r\neq0\), and interpret the singularity at the origin qualitatively.
\item A velocity field is \(\mathbf{v}=a\rho\erho+bz\khat\) in cylindrical coordinates. Find \(\nabla\cdot\mathbf{v}\). For what relation between \(a\) and \(b\) is the flow incompressible?
\item A purely azimuthal field is \(\mathbf{A}=A_\phi(r,\theta)\ephi\) in spherical coordinates. Write the general formula for its curl.
\item A tensor \(A\) has eigenvalues \(1,2,5\). Find \(I_1,I_2,I_3\) and write its characteristic polynomial.
\item For \(A_{ij}=x_{i}x_j\), compute \(\tr(A)\), \(A_{ij}A_{ij}\), and \(\partial_{j}A_{ij}\).
\item Use \(\varepsilon_{ijk}\varepsilon_{ilm}\) to derive Lagrange's identity \(|\mathbf{a}\times\mathbf{b}|^2=|\mathbf{a}|^2|\mathbf{b}|^2-{(\mathbf{a}\cdot\mathbf{b})}^2\).
\item Show that \(\nabla\cdot(f\mathbf{A}\times\mathbf{B})\) can be expanded using product identities.
\item In spherical coordinates, compute \(\nabla f\) and \(\nabla^2f\) for \(f=f(r)\), a radial-only scalar field.
\item In cylindrical coordinates, compute \(\nabla f\) and \(\nabla^2f\) for \(f=f(\rho)\), an axisymmetric radial field.
\item Explain why a formula that is true in Cartesian components may need correction terms in curvilinear coordinates if components are not physical components.
\item A complex number is multiplied by \(2\ee^{\ii\pi/3}\). Describe the geometric action.
\item Solve \(\ee^{2z}=1\) for complex \(z\).
\item Show that \(\cosh x\) and \(\sinh x\) parametrise the unit hyperbola \(X^2-Y^2=1\).
\item Compare \(\cos^2x+\sin^2x=1\) with \(\cosh^2x-\sinh^2x=1\). Explain the sign difference using exponentials.
\item Build a one-page personal formula sheet from memory covering all topics in this workbook.
\end{enumerate}
\end{practicebox}

\ifdefined\OTCOMBINED%
\else
\end{document}
\fi


\end{document}
